\subsection{The cohomology contribution to the division estimate}
\label{subsec:division-length}
\hypertarget{source-division-length}{}

The local cases bound lengths for individual characters. In the division lattice these lengths occur with eigenspace multiplicities and normalized place weights. We now sum them at one point of the auxiliary family at scale \(x\). Write \(z=x^\epsilon\),
\(b=e^z\), \(f=2n\) and \(B_0=fmbz\), with \(N_c\ge0\) the rank
budget for this point. Put \(\mathcal E=z(B_0+N_c)\).
\begin{definition}[Weighted local congruence numerator]\label{def:weighted-congruence-numerator}
Fix a lift \(1\le j\le m\). Let \(K\) be the coefficient field
and let \(\mu_0\) be the base tuple. The sum below ranges over
places \(w\) where the local tuple differs from \(\mu_0\).
Write \({\rm N}w\) for the norm of \(w\), \(e_w\) for its
absolute ramification index, and \(\eta_j(w)\ge0\) for the sum
of the normalized lengths at the local Hecke algebras.
Define the numerator \(H_j\) by
\[
 H_j=\sum_{w\ {\rm twisted}}
       \frac{\log {\rm N}w}{[K:\mathbb Q]} e_w\eta_j(w).
\]
For the coefficient DVR \(O_w\) at \(w\mid p\) and a finite-length module \(L\), normalization means
\[
 \operatorname{len}_w(L)=\frac{\operatorname{length}_{O_w}(L)}
 {\operatorname{length}_{O_w}(O_w/pO_w)}
 =\frac{\operatorname{length}_{O_w}(L)}{e_w}.
\]
\end{definition}
The normalized local congruence length and its behavior under
removal of eigenspaces admit the following direct description.

Let \(O\) be such a DVR and \(M\) the integral module at one index
over \(j,w\), viewed through its image in its generic fiber \(V\).
The finite Hecke algebra acts through an algebra of \(O\)-valued
functions on a finite set of characters and separates those characters.
For a set \(Z\) of good characters write \(V_Z=\bigoplus_{x\in Z}V_x\),
where \(V_x\) is the generic eigenspace, and write \(p_Z\) for the
corresponding rational projector.

\begin{lemma}[Return from the cut to the original Hecke lattice]
Let \(S\supseteq Z\) be the characters retained by the integral operator of Proposition~\ref{prop:cut-selector}. At a fixed
\(j,w\) let \(\eta_{\rm cut},\eta_{\rm original}\) be the normalized
congruence lengths for \(p_SM\) and \(M\), respectively. Then
\[
 0\le\eta_{\rm original}-\eta_{\rm cut}
           \le(\dim V_Z)v_p(h),
\]
where \(h\) is the finite row group order in
Proposition~\ref{prop:cut-selector}. If the total good multiplicity
at each twisted place is at most \(d_j\), the weighted increase in
the numerator is at most \(d_j\log h\).
\end{lemma}
\begin{proof}
The generic Hecke algebra contains the projector onto \(V_x\):
for each \(y\ne x\) choose \(a_y\) separating \(x,y\), and take the
product of \((a_y-a_y(y))/(a_y(x)-a_y(y))\). Hence
\(V=\bigoplus_x V_x\). For \(M\) the upper and lower lattices are
\[
 U=p_ZM=V_Z\cap(M+V_{Z^c}),\qquad B=M\cap V_Z.
\]
Both have full rank in \(V_Z\); at \(Z=\varnothing\) both are zero.

Now retain a set \(S\supset Z\) of eigenspaces and project the
integral module onto them. The integral lattice of the quotient
in its generic fiber is \(p_S M\). The upper lattice for \(Z\) there
is still \(U\) and its lower lattice becomes
\(B'=p_S M\cap V_Z\supset B\). The integral type operator
\(a\) of Proposition~\ref{prop:cut-selector} acts on \(V_Z\) by
the order \(h\) of the finite row group and vanishes on
\(V_{S^c}\).

For any \(v\in B'\), write \(v=p_S u\),
\(u\in M\). Then \(a u=h v\) lies in \(B\). Thus \(B'/B\) is
annihilated by \(h\) and generated by at most \(\dim V_Z\) elements,
since \(B'\) is a lattice over a DVR. Consequently the sequence
\[
 0\longrightarrow B'/B\longrightarrow U/B
        \longrightarrow U/B'\longrightarrow0
\]
computes the difference of the normalized lengths. Only the scalar
actions on \(V_Z\) and \(V_{S^c}\) enter: for \(p_Su=v\in V_Z\)
the other retained components of \(u\) vanish.

Suppose the total good multiplicity at each twisted \(w\) is at
most \(d_j\), the small-level rank at \(j\). For \(w\mid p\),
\(\log {\rm N}w=f_w\log p\) and
\(\sum_{w\mid p}e_w f_w=[K:\mathbb Q]\).
Consequently the total loss on returning to the uncut numerator is
at most \(d_j\sum_{p\mid h} v_p(h)\log p=d_j\log h\).
The admissible parameter bound \(\log h\le C_1 x^{C_1} z\)
is absorbed by \(C_1 fmbz^2\) beyond a single threshold, even
though the group may vary at every point. Thus an estimate on the
cut numerator with bracket
\[
 d_j\left\{C\big(fmbz^2+(\log x)^2N_c\big)+V\right\}
 \tag{*}\label{eq:local-division-star}
\]
gives the same form on the original numerator, with \(C\) increased
by \(C_1\).

\end{proof}
\paragraph{The residual local condition on the selected cut.}
Use Proposition~\ref{prop:weighted-local-return} on the good cut
of Proposition~\ref{prop:cut-selector} at each
point index above the family threshold. At a retained localization
let \(k\) be the residue field of \(O\), \(\bar\rho\) the
rank-\(r\) residual constituent, and \(\bar\epsilon\) the
residual cyclotomic character. Write \(S_i,Q_i\) here for the two
sets of places in the norm choice, with \(Q_i\) the auxiliary
twist level. Both lie away from \(\ell\).
\begin{lemma}[Local vanishing on the retained row]\label{lem:cut-residual-vanishing}
For each retained localization satisfying the cut conditions, every
\(v\in S_i\), and every integer \(e\), one has
\[
 H^a\big(E_v,\bar\rho\,\bar\chi_1^{-1}\bar\epsilon^{\,e}\big)=0
                       \qquad (a=0,1,2).
\]
\end{lemma}

\begin{proof}
\proofstep{Separate the chosen tame eigenvalue.}
At either type prime that occurs in \(S_i\), the projected
roots of the residual rank-\(r\) inertia polynomial avoid the
projected character label. Thus at the chosen tame generator
\(\gamma_v\), \(u=\bar\chi_1(\gamma_v)\) is not an eigenvalue
of \(\bar\rho(\gamma_v)\).
At the additional prime of \(S_i\), use the choice away from
\(\ell\) of the two distinguished order-\(15\) primes.

On the
full representation the inertia polynomial reduces at \(\gamma_v\)
to \((X-1)^r(X-u)^t\), by the simultaneous inertia decomposition
of Section~\ref{subsec:retained-galois}. The residual factorization
also gives the rank-\(r\) polynomial times \((X-u)^t\), since all
the cyclotomic twists agree on inertia. Cancelling that monic
factor identifies the rank-\(r\) polynomial with \((X-1)^r\).
The exclusion of the eigenvalue one also holds at \(\ell=3,5\), since
\[
 \Phi_{15}(1)=1\not\equiv0\pmod\ell.
\]
Thus \(u\ne1\). These arguments use the inertia representation at
the actual localization, so passage to its retained eigensystems
preserves the test.

\proofstep{Use local duality and the Euler characteristic.}
Cyclotomic inertia is trivial at \(v\). Hence \(\gamma_v-1\)
acting on \(M_e=\bar\rho\,\bar\chi_1^{-1}\bar\epsilon^{\,e}\)
is invertible, independently of \(e\). On its Tate dual the
inertia matrix is the inverse transpose, which has the same
fixed-space vanishing. For the dual over \(\mathbb F_\ell\)
one may identify the linear dual over \(k\) by the nondegenerate
finite-field trace pairing. Local Tate duality gives \(H^2=0\).
The Euler characteristic of a finite \(\ell\)-primary module
over this nonarchimedean local field of residue characteristic
different from \(\ell\) is one. Finiteness and \(H^0=H^2=0\)
therefore give \(H^1=0\). This proves the rank-\(r\) vanishing
at all the \(S_i\)-places used in the row presentations.
\cite{milne-duality}\label{cor:division-length-1}
\end{proof}

At \(Q_i\) the twist can have \(\ell\)-power inertia. The tame allowance at these places is the separate quantity
\[
 t\sum_{v\in Q_i}v(\zeta_v-1),
\]
as in Section~\ref{subsec:retained-twist}.
\cite{milne-duality}\label{cor:division-length-2}

The cut can be made at all indices of one family. Below the
threshold retain the whole space. Above it choose the good
subset in each localization as in Section~\ref{subsec:selected-cut}.
The quotient to the retained generic eigenspaces preserves the
point readings, the local twists and the block and weight labels.
It therefore carries each chosen auxiliary character to its
restriction on the surviving tests. The same character still
satisfies the simultaneous inequalities on these tests. The
level comparison and the group order in the return to the
original numerator are taken on that original point. In particular
the threshold for the restriction is independent of \(j,w\)
and of the translate. The constants in the local table on this
cut are chosen for the whole resulting family, before evaluating
it at the selected point. \label{cor:division-length-3}

\paragraph{Strict classes in the marked presentation.}
\begin{lemma}[Scalar strict contribution]\label{lem:scalar-strict-cost}
On the selected cut, let \(c_{\rm sc}\) be the contribution to the bracket in \eqref{eq:local-division-star} per unit of good multiplicity, from blocks \(2\le c\le t\) summed with the normalized coefficient-place weights. Then
\[
 c_{\rm sc}\le V\le\sqrt b.
\]
\end{lemma}

\begin{proof}
\proofstep{Identify the twists on the original field.}
In each local block the other conditions on the cocycle
are imposed before applying the strict budget. If \(\chi\) is the common auxiliary character on the two lines, the target-to-quotient ratio in block \(c\) is
\[
 (\chi\chi_\ell^{\,1-c})\chi^{-1}=\chi_\ell^{\,1-c}.
\]

The normalization is the same when the family is cut into the
retained eigenspaces. After the estimated local restriction maps
at the auxiliary primes and the row and column conditions have
been imposed, the scalar kernel for \(c\ge2\) is unramified
away from \(\ell\) and locally zero at every place dividing it.
Take \(a=c-1\), which runs through \(1,\ldots,t-1\).
The \(c=1\) kernel is bounded by the class number, and the
rank-\(r\) constituent uses its specialized signature Selmer
length. The same scalar ratios and kernels occur at both
type primes and in the complementary-prime calculation.

\proofstep{Normalize coefficient lengths.}
For the actual coefficient DVR \(O\), with fraction field \(F\),
there is a Galois-equivariant coefficient isomorphism
\[
 (F/O)(-a)
 \simeq O\otimes_{\mathbb Z_\ell}
             (\mathbb Q_\ell/\mathbb Z_\ell)(-a).
\]
The cochain complexes for these discrete coefficients commute
with this tensor product by choosing a finite
\(\mathbb Z_\ell\)-basis of \(O\). Kernels of localization
commute as well. A cyclic factor of length \(h\) over
\(\mathbb Z_\ell\) becomes \(O/\ell^h O\), of length
\(e_Oh\) over \(O\). Thus the normalized strict length of
the scalar kernel costs at most the corresponding summand
of \(V_\ell\) on this field. This also proves independence
of the choice of splitting coefficient extension.

\proofstep{Sum the strict bounds at the selected rank.}
The character ratios identify individual strict twists. Summation must stop below the selected rank on the original field, so comparison with the full interval estimate uses precisely its strict cutoff.

The bound for the weighted sum over primes from Section~\ref{sec:family-window} holds
for all finite subsets and its terms are nonnegative.
Taking their supremum gives
\[
 V=\sum_\ell(\log\ell)V_\ell\le\sqrt b.
\]

Each rational prime is charged after the sum of the good
multiplicities over localizations has been bounded by \(d_j\)
and after the coefficient place weights have been summed.
This puts \(V\) in the bracket of \eqref{eq:local-division-star} with coefficient one.
\label{cor:division-length-4}
\end{proof}

The small-level rank \(d_j\) equals the base row
dimension \(D_0\) by Proposition~\ref{prop:base-theta-rank}. The base-level theta
map surjects onto the base row space and its adjoint has
rank \(D_0\). The integral projection denominators are
included in the length comparison above.

\subsubsection*{Summing the local bounds}
\hypertarget{source-numerator-sums}{}

A bound for a congruence ideal must first be turned into a bound for the corresponding integral module. This conversion keeps the rank of the retained eigenspace in the sum over good characters.

\begin{established}[Hecke congruence ideal and normalized length]\label{hyp:eigencharacter-lattices}
First consider the passage from a Hecke congruence ideal
to the congruence length of its module. Work over one
coefficient DVR \(O\), normalize \(v(\ell)=1\), and write
\(Y\) for the theta eigensystems and \(Z\subseteq Y\) for
the good ones. Let \(T\) be the Hecke algebra inside the
algebra of functions on the generic eigensystems. Write \(I_Y=\{f\in T:f|_{Y^c}=0\}\). For each \(\pi\in Z\) its evaluation and normalized colength satisfy
\[
 0\ne I_Y(\pi)=(f_\pi(\pi))\subseteq O,\qquad
 v_{cc}(\pi)=v(f_\pi(\pi)),\qquad f_\pi\in I_Y .
\]
\end{established}
\begin{proof}
For each generic character \(y\ne\pi\), choose \(t_y\in T\) with \(t_y(\pi)\ne t_y(y)\). The product \(\prod_{y\ne\pi}(t_y-t_y(y))\) vanishes away from \(\pi\) and has nonzero evaluation there. The scalars \(t_y(y)\) belong to the splitting DVR \(O\). Hence the image of the supported ideal is a nonzero ideal of \(O\). Choose a generator and lift it to a supported element of \(T\).
\end{proof}

\begin{lemma}[From eigencharacter valuations to module colength]\label{lem:eigencharacter-colength}
With the ideal and length of Lemma~\ref{hyp:eigencharacter-lattices},
let \(M\) be the Hecke-stable lattice and \(V_Z\) its good generic
subspace. Put \(U=p_ZM\), \(B=M\cap V_Z\)
and let \(a\) be the integral type operator of
Proposition~\ref{prop:cut-selector}. Then the normalized
colength of \(B\) in \(U\) is at most
\[
 \sum_{\pi\in Z}\dim(V_\pi)\bigl(v_{cc}(\pi)+v_\ell(h)\bigr).
\]
\end{lemma}

\begin{proof}
The type operator \(a\) acts by \(h\) on the good
eigenspaces and by zero on the remaining eigenspaces
of \(Y\). The product \(af_\pi\) kills every eigenspace
outside \(Z\). It also preserves the integral module,
so it maps \(U=p_ZM\) into \(B=M\cap V_Z\), by applying
it to a lift in \(M\). The products for the different
\(\pi\) may all act on multiple good eigenspaces.
Order those eigenspaces and filter \(U\) and \(B\) by
intersection with their successive partial sums.

For
the step adding \(V_\pi\), the graded lattice of \(U\)
is the projection of the intersection at this step
onto \(V_\pi\), and likewise for \(B\). Multiplication
by \(h f_\pi(\pi)\) takes the first graded lattice into
the second, since \(af_\pi\) preserves the partial sum.
For the step adding \(V_\pi\), the DVR index formula therefore bounds the normalized graded length by
\[
 \dim V_\pi\bigl(v_{cc}(\pi)+v_\ell(h)\bigr).
\]
With \(U_i,B_i\) the intersections with the first \(i\) eigenspaces,
the successive index is read in the projection to the new eigenspace.
Equivalently
\[
 \operatorname{len}(U/B)=
 \sum_i\operatorname{len}\bigl((U_i/U_{i-1})/(B_i/B_{i-1})\bigr)
 \le\sum_{\pi}\dim V_\pi\bigl(v_{cc}(\pi)+v_\ell(h)\bigr).
\]
This proves the local comparison, with the Hecke
algebra and the theta set of the cut when it is
applied to the cut.
\end{proof}

\begin{proposition}[Weighted local return on the cut]\label{prop:weighted-local-return}
Fix a lift on the selected point family and use its good
multiplicities from Proposition~\ref{prop:cut-multiplicity}.
The numerator of Definition~\ref{def:weighted-congruence-numerator} satisfies \eqref{eq:local-division-star} of Section~\ref{subsec:division-length} on the cut, with the allowed additional term \((\log x)^2N_c\) and a threshold uniform over the selected points. In the notation of the local table its contribution at an ordinary localization \(i\), away from the two type primes, satisfies
\[
 \eta_i\le \sum_{\pi\in Z_i^+}m_{i,\pi}
               \bigl(v_{cc}(i,\pi)+v_\ell(h)\bigr).
\]
Here \(v_{cc}(i,\pi)\) is the local congruence bound, \(h\) the order of the finite row group and \(m_{i,\pi}\) the good multiplicity. The separate loss \(d_j\log h\) from returning to the original Hecke module is absorbed at this scale.
\end{proposition}

\begin{proof}
\proofstep{Weight each localization by its normalized local degree and good multiplicities.}
Apply Lemma~\ref{lem:eigencharacter-colength} on the cut at
a single selected point. Fix a lift \(j\) and write \(d_j\) for the rank after
the small-level rank identification. Let \(\mathcal T\)
be the finite set of rational primes below the twisted
coefficient places in its support. The type and auxiliary prime cases use
Propositions~\ref{prop:type-prime-row-length}
and~\ref{prop:auxiliary-prime-row-length}. The local length
comparison at a prime \(\ell\) other than
\(g_0,g_1\), is the bound
\[
 \eta_i\le \sum_{\pi\in Z_i^+}
      m_{i,\pi}\bigl(v_{cc}(i,\pi)+v_\ell(h)\bigr).
\]
Here \(\eta_i\) is the normalized congruence length of the
cut localization, \(m_{i,\pi}\) is its multiplicity, and
\(v_{cc}\) is the local congruence bound attached to the
Galois family at \(\pi\). The factor \(h\) belongs to this
comparison on the cut itself. Returning to the original
module incurs the separate loss proved above.

Lemma~\ref{lem:marked-row-column-lattices} gives the required integral presentation at each good character of this cut. Proposition~\ref{prop:marked-signature} supplies its complement estimate. The resulting local table has
\[
 v_{cc}(i,\pi)\le
 \nu_\ell+2\ell^{-2}+\mathcal L_\ell+V_\ell+v_\ell(h_E),
\]
where \(\nu_\ell\) and \(\mathcal L_\ell\) are independent
of \(i,\pi\) for fixed \(\ell\) at this point. The table
incorporates the uniform bound on \(s_\pi\) separately
from the split row estimate. At each twisted \(w\mid\ell\), \eqref{eq:retained-multiplicity}
gives
\[
 \eta_{j,w}=\sum_{i\text{ over }(j,w)}\eta_i
 \le d_j\bigl(\nu_\ell+2\ell^{-2}+\mathcal L_\ell+V_\ell
                 +v_\ell(h_E)+v_\ell(h)\bigr).
\]
All terms in this upper bound are nonnegative and independent of
\(w\). Multiply by \(e_w\log\mathrm Nw/[K:\mathbb Q]\) and sum
over the chosen places. The factor from their weights is at most
\(\log\ell\), by the residue-degree formula on all places above
\(\ell\). At the two type primes the
paired-character estimate bounds the weighted lengths
directly by
\(d_j\log\ell(\mathcal L_\ell+2\ell^{-2}+V_\ell+v_\ell(h_E))\).
Up to nonnegative terms the same summation applies there.

The class number term contributes at most \(\log h_E\),
since
\[
 \sum_{\ell\in\mathcal T}v_\ell(h_E)\log\ell=\log\!\left(\prod_{\ell\in\mathcal T}\ell^{v_\ell(h_E)}\right)\le\log h_E.
\]
The factor \(h\) costs at most \(\log h\)
by the same argument. Also
\(\sum_{\ell\in\mathcal T}\ell^{-2}\log\ell\) is absolutely
bounded by the series over the integers. The strict
cohomology terms contribute at most the global budget
\[
 V=\sum_\ell V_\ell\log\ell
\]
once, since \(\ell\) is
summed only after the multiplicity and degree estimates.
This bound is independent of the number of localizations
and the coefficient-field degree.

The scalar cyclotomic term has its own bound for the weighted sum over primes.
Put \(a=tr\), which is odd since \(t\) is odd, \(k,m\)
are even, and \(r=k+m-1\). The cyclotomic minimum \(\nu_\ell\) in \eqref{eq:signature-cyclotomic-minimum}
is taken over the absolute Galois group of this same field \(E\). If \(\ell\mid Q_S\), the image
contains \(1+\ell^{v_\ell(Q_S)}\) since \(E\) lies in
the cyclotomic field of conductor \(Q_Sqd\). The
binomial argument in Section~\ref{subsec:split-local-length}
gives
\[
 \nu_\ell\le v_\ell(Q_S)+v_\ell(a).
\]

The depth at \(2\) is at least two. If
\(\ell\nmid Q_Sqd\), the cyclotomic image is full by
the Chinese remainder theorem. Then \(\ell\) is odd
and \(-1\) gives \(\nu_\ell=0\).
At \(q,d\), the host Artin action of the factor at
that prime has order dividing two. Thus the cyclotomic
image modulo that prime contains the subgroup of
squares, of order \((\ell-1)/2\). This order exceeds
\(a\) by the auxiliary prime choice, so a generator raised
to \(a\) differs from \(1\). This again gives
\(\nu_\ell=0\). Therefore
\[
 \sum_{\ell\in\mathcal T}\nu_\ell\log\ell
 \le \log Q_S + \log(tr)=O(\log x).
\]
Here \(Q_S\le4 f_\psi n\le n^2\) eventually by
\eqref{eq:host-character-compositum} and \eqref{eq:dyadic-coordinate-orders}.
The target is fixed before the scale and the ranks grow polynomially.

The multiplicity estimate and the local table join at
one common threshold in \(x\). The arguments of
Section~\ref{subsec:special-prime-lengths}
give the type and ramified entries. Section~\ref{subsec:ordinary-split-length}
proves the ordinary split entry and the implication from a local
inert band to the indicator in the sum. The cyclic row estimates
in Section~\ref{subsec:inert-rows} supply the long-orbit and
the two remaining inert entries. The comparison above the
cutoff in Section~\ref{subsec:outside-windows} uses the
integral local classification on the same jump windows. The table
assigns \(\mathcal L_\ell\) in the following ranges.

\begin{itemize}
\item At each type prime \(g_i\) the table entry and sizes satisfy
\[
 \mathcal L_{g_i}\le Cx^C\eta_c,\qquad \log\eta_c\le x^{\epsilon/4},\qquad \log g_i=O(z).
\]

\item At each dyadic or odd conductor prime and at three,
Proposition~\ref{prop:small-split-loss} bounds the cost by \(C fmb\log x+Cx^Cz^2\). The weighted
sum is \(O(fmbz^2)\) by the size comparison of
Section~\ref{subsec:split-local-length}.
\item At \(q,d\), the cost is at most \(2C fmb\).
The weighted sum is \(O(fmbz)\) by the conductor of the auxiliary field
bound on the chosen prime pair.
\item The remaining table uses a positive cutoff
\(\ell_*\). The summation below gives \(O(fmbz^2)\)
for the boundary, split, inert and short-band costs
weighted by \(\log\ell\) at primes at most \(\ell_*\).
\item Above \(\ell_*\) the remainder is supported on
the exceptional set. With
\[
 U=x^{C+2}\exp(x^{\epsilon/4}),
\]
its cardinality is at most \(388U^8\). The weighted cost
per exceptional prime is at most \(3000(C+1)^2U^{10}\),
giving at most \(1164000(C+1)^2U^{18}\) in total.
\end{itemize}

\proofstep{Sum the same table before absorbing the exceptional-prime cost.}
The table is now weighted by the local degrees and actual multiplicities. Summing those entries first keeps the exceptional primes and the scalar cohomology on the same retained family.

For fixed \(C\) and \(\epsilon>0\),
\[
 \log(U^{18}/b)=18(C+2)\log x+18x^{\epsilon/4}-x^\epsilon
 \longrightarrow -\infty .
\]
In particular the exceptional sum and the two type-prime
costs are \(o(b)\) uniformly. The same comparison absorbs
any fixed polynomial in \(x\). The bound on \(\log h_E\)
can be obtained separately. All character conductors of the abelian
host divide \(N\), so the conductor--discriminant formula
gives \(\log|\Delta_E|\le [E:\mathbb Q]\log N\).
Minkowski's bound supplies an ideal in every class of
norm at most \(2^{[E:\mathbb Q]}|\Delta_E|^{1/2}\).
\cite[Theorems 3.27 and 4.3 respectively]{milne-cft,milne-ant}
With \(n_E=[E:\mathbb Q]\) and
\(D=2^{n_E}|\Delta_E|^{1/2}\), count these representatives by
\[
 h_E\le\sum_{\mathrm N\mathfrak a\le D}1
  \le D^2\sum_{\mathfrak a}(\mathrm N\mathfrak a)^{-2}
  =D^2\zeta_E(2)\le D^2\zeta(2)^{n_E}.
\]
The last inequality is Euler factor by Euler factor, since there
are at most \(n_E\) prime ideals above each rational prime and
their residue degrees are at least one.

Taking logarithms and
using \(\log|\Delta_E|\le n_E\log N\) gives
\(\log h_E=O(n_E(1+\log N))\), which is polynomial in \(x\)
under the selected size bounds. With the local table,
these estimates give the cut numerator bound (*) with the
allowed additional
nonnegative term \((\log x)^2 N_c\). The loss \(d_j\log h\)
from the cut to the original level can be absorbed
at this same scale.

The local inputs are assembled on the cut family
before applying the weighted sum. The estimate uses
the complete local table and the rank and period
identification at the chosen points. \label{cor:numerator-sums-1}
\end{proof}

\subsubsection*{Summation below the prime cutoff}
\label{subsec:prime-cost-sums}
\hypertarget{source-prime-cost-sums}{}

The remaining sum below the cutoff contains contributions from both split and inert places as well as boundary terms. Their individual bounds have to be summed with the prime weights before the total can be compared to the allowed numerator. Write \(z=x^\epsilon\), \(b=\exp z\), and
\(H=b+S_c+\eta_c+t+2\). Let \(k\) be the ambient rank parameter,
and put \(N(m)=k+m(t+1)\).
Write \(\delta\) for the density of the selected indices with the larger weights.
\begin{established}[Size inequalities for the local table]\label{hyp:small-prime-table-ranges}
The cutoff and the secondary threshold are
\[
 \ell_*=\frac{mb}{n t^2(N(m)+t)z^2+k n^2(S_c+t+z)},
 \qquad A=\frac{kz}{m\delta}.
\]
Write \(X'=\exp(x^{\epsilon/3})\), \(b'=xb/X'\), and
\(\omega=b/(nX')\). Let \(\eta_3=\eta_c+t+2\), let \(w_P\) be
the upper place-weight size plus \((t+1)/2\), and put
\(G_P=w_P+S_c+\eta_3\). For the selected ranks and weights, beyond a threshold depending on the fixed parameters,
\[
  \ell_*\le b/z,\quad z\le A,\quad r/A\le 2m/z,\quad(3+\log A)^2\le z,\quad H\le2b,\quad z\le b'.
\]
\[
 \begin{gathered}
 \eta_3\le4\eta_c,\qquad w_P+\eta_3\le9b',\qquad G_P\le10b',\\
 r b'A\log x\le mb,\qquad r\eta_3\le4m\omega,\qquad
 C(\eta_3+S_c+z)\le\omega/8 .
 \end{gathered}
\]
\end{established}
\begin{proof}
Lemma~\ref{lem:window-parameter-choice} fixes the ranks before the scale and the weights of the retained positions before choosing the place. In this
calculation \(x^{-\epsilon}\le\delta\le1\), \(r\le2k\),
\(R\le N(m)+t\), \(m\le k\), and the ranks are at most a fixed
power of \(x\). In particular \(\ell_*\le b/z\), and
\[
 z\le A,\qquad r/A\le 2m/z,\qquad
 (3+\log A)^2\le z
\]
beyond one threshold. The last inequality follows from
\(A\le k z^2\) and \(z/(\log x)^2\to\infty\).
The size inputs also give \(H\le2b\).

At the smaller places the two possible centers satisfy
\[
 \omega\le 5b'/4,\qquad b_v+\lfloor b'/x\rfloor\le5b'/4+b'/x .
\]
Here the bound for \(b_v\) is from Section~\ref{subsec:exceptional-range}. The widths \(S_c,\eta_c,t+2\) are \(o(b')\), and \(t+2\le3\eta_c\). Including these widths and the half shift gives \(w_P+\eta_3\le9b'\) and \(G_P\le10b'\) uniformly.

For the fourth inequality the selected reserve bounds
\(r(b'+S_c)A\log x\) by \(\delta mb\). The fifth follows from
the reserve \(N(m)\eta_c\le\delta m\omega\) and \(r\le N(m)\).
For the last inequality, the sizes of \(\eta_c,S_c,z\) are
\(O(\omega/\sqrt x)\), with constants fixed before taking \(x\)
large. For fixed \(\epsilon>0\) this follows also by comparing
their bounds, \(x^{O(1)}\exp(x^{\epsilon/4})\), with \(\omega\).
We have \(b'\ge z\) at large \(x\).
\end{proof}

\begin{proposition}[Summation up to the small-prime cutoff]\label{prop:small-prime-cost-total}
Under the cutoff and size inequalities of Lemma~\ref{hyp:small-prime-table-ranges},
let \(c_\ell\) be the sum of the local boundary, split and inert charges in the table at \(\ell\), before multiplication by \(\log\ell\). Then
\[
 \sum_{\ell\in P}(\log\ell)c_\ell\ll_C fmbz^2\qquad\bigl(P\subseteq\{\ell:5\le\ell\le\ell_*\}\text{ finite}\bigr).
\] The threshold
is uniform in the retained point and the subset of primes.
\end{proposition}

\begin{proof}
For a set of integers \(2\le a\le L\), summing over all such
integers gives
\[
 \sum \log a\le L\log L,\qquad
 \sum \frac{\log a}{a}\le (1+\log L)^2.
\]
The second bound uses the harmonic sum and \(\log a\le\log L\).
Replacing \(a\) by \(a-1\) in the denominator costs at most a
factor two. These bounds apply to any subset of the indicated primes.
The boundary cost \(Cfn t^2 R z^2\) consequently contributes at
most \(C fmbz\), by the first term of the cutoff denominator.
\proofstep{Split places.}
The split cost is \((2fm+f)e_P(\ell)\), where
\(e_P(\ell)=C(1+H/\ell+\log_\ell H)\). The other uses of \(e_P\)
have coefficient at most \(2Cfm\). After multiplication by
\(\log\ell\), the three terms in \(e_P\) sum to
\[
 O_C\!\left(\ell_*\log\ell_*+
             H(1+\log\ell_*)^2+\ell_*\log H\right)
       =O_C(bz^2).
\]

\proofstep{Inert places below the secondary threshold.}
For inert \(\ell<A\) of order exceeding \(4L_w\) in the
host's real cyclotomic quotient, the cost is \(4Cfr\) times
\[
 J_\ell=\frac{H}{(\ell-1)^{j_\ell+1}}
         +(G_P+1)(j_\ell+1)+z,\qquad
 j_\ell=\max(1,\lceil\log_{\ell-1}A\rceil).
\]
For \(\ell\ge5\), \((\ell-1)^{j_\ell}\ge A\) and
\(j_\ell+1\le3+\log A\). The logarithm-weighted sum of the
first term, including \(r\), is
\(O(rH(1+\log A)^2/A)=O(mb)\).
For the other terms it is
\(O(rb'A(3+\log A)^2)=O(mbz)\). This uses the reserve for
\(rb'A\) and \(G_P+1=O(b')\).

The other inert primes have a contribution \(4CfrH/\ell^2\).
At \(\ell\ge A\) use
\[
 \sum_{a\ge A}(\log a)/a^2=O((1+\log A)/A)
\]
by integral comparison, and \(r/A\le2m/z\). At \(\ell<A\)
the order is at most \(4L_w\); here the ordinary tail bound at
\(A\) gives no control. Put \(R_0=n/z\). The parameters of the auxiliary construction
give \(3\le R_0\), \(k/m\le R_0^2\), \(n\le o2^e\) and
\(2oL_w\le z\), where \(e\) is the host's dyadic exponent
and \(o\) its fixed cyclic order parameter. Thus
\[
 2^e/L_w\ge2R_0.
\]

The quotient
\((\mathbb Z/2^{e+2}\mathbb Z)^\times/\{\pm1\}\) has order \(2^e\),
so every order \(h\) on it is a power of two. By
Lemma~\ref{lem:short-real-dyadic-order} an integer of order \(h\) is
\(\pm1\pmod{2^{e+2-v_2(h)}}\).

Consequently \(h\le4L_w\) forces the original integer to
be \(\pm1\bmod M\), where
\[
 M=2^{e+2-\lfloor\log_2(4L_w)\rfloor}
       \ \ge\ 2^e/L_w\ \ge\ 2R_0 .
\]
The same congruence holds if \(a=1\). There is no term at
\(1\) in the sum of interest. All the other positive integers
in these two residue classes are \(jM\pm1\) with \(j\ge1\).
They are between \(jM/2\) and \(2jM\), so
\[
 \sum_{\substack{a\ge2\\a\equiv \pm1\bmod M}}
            \frac{\log a}{a^2}
   \le \frac8{M^2}\sum_{j\ge1}\frac{\log(2jM)}{j^2}
   =O\!\left(\frac{1+\log M}{M^2}\right).
\]
The function \((1+\log u)/u^2\) is decreasing for \(u\ge1\).
Multiplication by \(r\le2mR_0^2\) therefore gives
\(O(m(1+\log R_0))\). This bound accounts for the small
inert primes using only the dyadic quotient, with no estimate
for primes in arithmetic progressions.

The remaining inert terms cost \(4Cfr\) times
\[
  \frac{w_P+\eta_3}{\ell}+S_c+z+\eta_3 I_\ell .
\]
Here \(I_\ell\) is the indicator that some \(s\) between
\(\omega\) and the upper place-weight size before the
\((t+1)/2\) shift satisfies
\[
 |\ell s-b|\le
   C\bigl(\ell(\eta_3+S_c+z)+b/\ell+s\bigr).
\]
For \(\ell\ge2C\) the right side is at most
\(\ell s/8+b/2+\ell s/2\), by the selected reserves.
It follows that \(\omega\le s\le4b/\ell\) in this band.
Hence \(r\eta_3 I_\ell\le16mb/\ell\).

At the finitely
many \(\ell<2C\) we use \(r\eta_3\le4m\omega\le4mb\).
Also \(r(w_P+\eta_3)\le9rb'\le9mb\), and the second
term of the cutoff denominator gives
\[
 \ell_*r(S_c+z)\le2mb.
\]
The first and band terms
are summed using \(\sum (\log\ell)/\ell\), and the
constant term using \(\sum\log\ell\). Their total
contribution is \(O_C(fmbz^2)\).

The estimates above apply together to any finite subset of
the primes \(5\le\ell\le\ell_*\). They give the desired
\(O_C(fmbz^2)\) bound for the numerical table used in the
weighted export. The threshold is common to all the retained points with the specified size inputs.
The sum is therefore available for the actual row once each localization satisfies the corresponding local Galois estimate. \label{cor:prime-cost-sums-1}
\end{proof}

\subsubsection*{The exceptional range}
\label{subsec:exceptional-range}

Above \(\ell_*\), only the exceptional jump windows require this additional sum. We count the primes they can contain and bound the total weighted cost uniformly in the window data. Write
\[
 \mathcal D=mb/\ell_*,\qquad
 U=x^{C+2}\exp(x^{\epsilon/4}),\qquad
 s_+=\lfloor2b/\ell_*\rfloor .
\]
The selected parameters give \(U\ge1\),
\(n,k,t,m,z,\eta_c\le U\), \(f,r\le2U\),
\(R,N(m)+t\le4U^2\) and \(S_c\le4U\).
Fixed polynomial rank bounds suffice after taking
\(C\) large enough. The two summands of \(\mathcal D\)
are at most \(4U^7\) and \(6U^4\) respectively, hence
\[
 \mathcal D\le10U^7.
\]
Since \(U^{14}=o(b)\), this
gives \(\ell_*^2>8b\) beyond a common threshold.
Also \(\ell_*>A\), since \(A\le k z^2\) grows
polynomially in \(x\).

The centers \(W\) of the exceptional windows are
\(b,\omega\), and \(b_v+\lfloor b'/x\rfloor\) for the
top places \(v\). Their number is at most \(2n+2\),
using the bound \(2n\) for the number of top places.
The weight construction gives
\(\omega\le b_v\le 5b'/4\), so each center lies in
\([\omega,2b]\) for large \(x\). The window about \(W\)
contains the integers in
\[
[W-\eta_c-2s_+-3,\,
  W+S_c+2\eta_c+2s_++4].
\]
All these integers are positive and less than \(3b\).
Indeed \(s_+\le2\mathcal D/m\) and
\(U^7=o(\omega)\) since \(n\) is polynomial,
\(\omega=b/(nX')\) and \(\log X'=x^{\epsilon/3}\).
This controls both endpoint displacements.

The exceptional set consists of primes \(\ell>\ell_*\)
dividing one of these integers. Two distinct such
divisors would have product greater than
\(\ell_*^2>8b\), so each integer contributes at most
one prime. An interval of length \(D\) contains at
most \(D+1\) integers. Thus the set has size at most
\[
 (2n+2)(S_c+3\eta_c+4s_++8)
 \le(4U)(97U^7)=388U^8.
\]
Each prime in the set is less than \(3b\), so
\(\log\ell\le3U\).
\label{cor:prime-cost-sums-2}

For these primes the boundary cost is
\(Cfn t^2 R z^2\le8C U^8\). The envelopes for the
pointwise row costs are the split cost and the second
inert cost, since \(\ell>\ell_*>A\). For \(\ell>7\),
\[
\begin{split}
 H/\ell^2 &\le H/\ell \le 2\mathcal D\le20U^7,\\
 e_P(\ell)&=C(1+H/\ell+\log_\ell H)\le23C U^7.
\end{split}
\]
For the flat part of the inert cost, the numerical
expression is
\((w_P+\eta_3)/\ell+S_c+z+\eta_3 I_\ell\),
where \(I_\ell\in\{0,1\}\) is the band indicator.
Here \(b'<\ell_*\) at the common threshold, by
\(\mathcal D x/X'\to0\), and \(w_P+\eta_3\le9b'\).
Thus the flat part is at most
\(9+4U+U+4U\le18U^7\).

The split envelope is at most \(138C U^9\).
The inert envelope is at most
\(608C U^9+92C^2U^9\), using
\[
 4Cfr(20+18)U^7+2Cfm(23C U^7)
       \le (608C+92C^2)U^9.
\]
Adding the boundary cost and multiplying by
\(\log\ell\) gives the bound
\(3000(C+1)^2 U^{10}\) per prime.
This proves the exceptional contribution to the
weighted sum from the pointwise row comparisons.
The complementary large-prime range uses the
congruence estimate with zero additional length
outside these windows.
\label{cor:prime-cost-sums-3}



\begin{proposition}[The numerator on the original lift]\label{prop:original-weighted-numerator}
Use the completed family with local contributions from
Lemma~\ref{lem:cut-residual-vanishing},
Lemma~\ref{lem:scalar-strict-cost} and the summation in
Section~\ref{subsec:division-length}. Include the correction
for the return from the cut in \(C\). Then beyond a common
threshold, at every point of the family and for every lift \(j\), the
weighted numerator on the original module satisfies
\[
           H_j\le D_0(|C|+1)z(B_0+N_c).
\]
\end{proposition}

\begin{proof}
Apply \eqref{eq:local-division-star} on the original module. The rank equality and the
inequalities \((\log x)^2\le z\), \(z\ge1\), and \(V\le\sqrt b\le B_0\)
give the stated bound. Choose the maximum of the thresholds for
Lemma~\ref{lem:scalar-strict-cost}, the cut comparison,
Proposition~\ref{prop:base-theta-rank} and the logarithm estimate.
Each threshold applies to the family before choosing its point or
lift.
\label{cor:division-length-5}
\end{proof}

The bound controls the cohomological term in the dual-lattice index
below. The first division error uses the transported row indices and
Weyl inclusion on the actual integral lattices of the same member in
Section~\ref{sec:word-division}. Both estimates and the period identities
are used in \eqref{eq:division-slopes} with the fixed exponent \(e_D=2\).

\paragraph{The contribution of the numerator to the dual index.}
\begin{established}[Nested lattices at the twisted place]\label{hyp:signed-dual-lattices}
At the used tuple, let \(\Lambda_\Delta\subseteq\Lambda_\nabla\)
be the smaller and ordinary section lattices of the lift in the
whole block. Transfer the two into the same coordinate space as in
Section~\ref{subsec:selected-period}, and put
\[
 \mathcal T=P_j(E^{-1}\Lambda_\nabla),\qquad
 \mathcal C=(VZ)^{-1}\Lambda_\nabla,\qquad
 \mathcal D=(VZ)^{-1}\Lambda_\Delta .
\]
These full lattices satisfy \(\mathcal D\subseteq\mathcal C\subseteq\mathcal T\). In the notation at the beginning of the section,
\(e_w\eta_j(w)=\operatorname{length}(\mathcal T_w/\mathcal C_w)\)
at the twisted place. Let \(J_w\) be the least nonnegative
integer such that \(\mathfrak p_w^{J_w}\mathcal C_w\subseteq
\mathcal D_w\). The congruence length uses \(\mathcal C_w\)
inside \(\mathcal T_w\). The exponent \(J_w\) measures
the further inclusion of \(\mathcal D_w\) into \(\mathcal C_w\).
\end{established}
\begin{proof}
The transfer identifies the generic spaces, and the block projection is surjective, so the three images span their prescribed spaces. The first inclusion follows by inverse image from \(\Lambda_\Delta\subseteq\Lambda_\nabla\). If \(v\in\mathcal C\), then \(EI_jv=VZv\in\Lambda_\nabla\). Thus \(I_jv\in E^{-1}\Lambda_\nabla\), and applying \(P_jI_j=1\) shows that \(v\in\mathcal T\). Localization preserves both inclusions. The quotient of the second pair has the stated normalized length by the definition of \(\eta_j(w)\), and a power of the maximal ideal kills the finite quotient of the first pair.
\end{proof}

\begin{proposition}[The local signed dual index on the common lift]
\label{prop:signed-lift-dual-index}
For the three lattices and \(J_w\) of
Lemma~\ref{hyp:signed-dual-lattices}, let \(D_G\) clear
the positive class and level denominators. Suppose \(s\) is a
positive integer satisfying the classwise value dual inclusion.
Then
\[
 [\mathcal T_w:(\mathcal T^\#)_w]_w
 \le e_w\eta_j(w)+D_0\bigl(J_w+v_w(sD_G)\bigr).
\]
If the same integer \(s\) works at every place in the sum and
\(D_*\) is a common multiple of the \(D_G\) at that point, divide
by \(D_0\) and sum with the normalized place weights. The
denominator correction is then at most \(\log D_*+\log s\).
The least exponent \(J_w\) is charged to the first error.
\end{proposition}

\begin{proof}
\proofstep{The class denominators and the global Hermitian dual.}
The Hermitian form on these coordinates is the restriction of the
level form. It is a sum of the class pairings divided by
the positive class weights, with the level normalization included.
Let the positive integer \(D_G\) be a common multiple of
these denominators, and suppose the positive integer \(s\)
satisfies the value-lattice
dual inclusion at every class. This means that the class pairing
of an ordinary value with \(s\) times any smaller-lattice
value is integral.

For class weight \(d_c\) and local pairing \(B_c\), the contribution of ordinary \(u\) and smaller-lattice \(v\) is integral after scaling:
\[
 sD_G\,\frac{B_c(u,v)}{d_c}
     =\frac{D_G}{d_c}B_c(u,sv).
\]
Here \(d_c\mid D_G\) and the value-lattice hypothesis makes \(B_c(u,sv)\) integral. Summing puts \(sD_G\) times the smaller section lattice in the ordinary dual. Restriction to the level subspace
uses the same class pairing and its normalizing denominator.
For \(z\in\mathcal D\) and \(y\in E^{-1}\Lambda_\nabla\)
orthogonality gives
\[
 \langle z,P_jy\rangle=\langle I_jz,y\rangle
\]
for these restricted forms. It follows that
\(sD_G\mathcal D\subseteq\mathcal T^\#\), where
\(\mathcal T^\#\) is the Hermitian dual over \(\mathcal O_K\).
Therefore
\[
 \mathfrak p_w^{J_w+v_w(sD_G)}\mathcal C_w \subseteq(\mathcal T^\#)_w.
\]
This inclusion concerns the
localization of the global dual and remains valid at places
interchanged by the coefficient involution.

\proofstep{The signed determinant index at the chosen place.}
For full local lattices take the signed index \([L:L']_w\)
to be the valuation of the determinant of a basis of \(L'\)
in a basis of \(L\). These indices add in triples of lattices
and increase as \(L'\) decreases. With \(d_j(w)=D_0\) the last
inclusion gives
\[
 [\mathcal T_w:(\mathcal T^\#)_w]_w
 \le
  e_w\eta_j(w)+D_0\bigl(J_w+v_w(sD_G)\bigr).
\]
\proofstep{Normalization of the place sum.}
Put \(\omega_w=\log\mathrm Nw/[K:\Q]\). Dividing the
local inequality by \(D_0\) and summing at the twisted places
gives
\[
 D_0^{-1}\sum_{w\ {\rm twisted}}\omega_w
        [\mathcal T_w:(\mathcal T^\#)_w]_w
 \le H_j/D_0
     +\sum_{w\ {\rm twisted}}\omega_w J_w
     +\sum_{w\ {\rm twisted}}\omega_w v_w(sD_G).
\]
Here \(H_j=\sum_{w\ {\rm twisted}}\omega_w e_w\eta_j(w)\).
Choose \(D_*\) divisible by every denominator \(D_G\) at
this point, and use the same \(s\) at every place. Then
\(0\le v_w(sD_G)\le v_w(sD_*)\). For a positive integer \(a\),
\[
 \sum_{w\mid p}\omega_w v_w(a)
   =\frac{\sum_{w\mid p}e_w f_w}{[K:\Q]}v_p(a)\log p
   =v_p(a)\log p,
\]
where \(f_w\) is the residue degree over \(p\).
This gives the upper bound \(\log D_*+\log s\) on the
last sum, also for a partial set of places. Thus the weighted signed index per lift,
divided by the positive row rank of the selected bundle,
is at most
\[
        (|C|+1)z(B_0+N_c)+\log D_*+\log s
        +\sum_{w\ {\rm twisted}}\frac{\log\mathrm Nw}{[K:\Q]}J_w .
\]

In the
global slope comparison the term \(J_w\) is charged in
the same per-lift allowance as the row comparison error.
Write \(J_{j,w}\) for this least exponent and \(e_w\) for
the ramification index. The chosen nonnegative errors
\(r_{j,w}\) bound the paired row-index remainder after
summing over lifts and including the declared \(2/p^2\)
term. The per-lift allowance \(E_{1,j}\)
must satisfy
\[
  \sum_{w\ {\rm twisted}}\frac{\log\mathrm Nw}{[K:\Q]}
                  (e_w r_{j,w}+J_{j,w})\le E_{1,j}.
\]
In the numerical scale the normalized exponent
\(J_0(j,w)\) is \(J_{j,w}/e_w\). Separately, the
allowance from the numerator divided by \(D_0\), plus
\(\log D_*+\log s\), is the quantity denoted \(E_{2,j}\)
in the division estimate. Thus the preceding signed-index
estimate gives the two terms with their specified
normalizations.

\proofstep{The rational-prime majorant for the least exponent.}
For the integral inclusion of lift lattices, let \(J(j,p)\) be the rational-prime exponent that bounds the local exponents:
\[
                     J_{j,w}\le e_w J(j,p)\qquad(w\mid p).
\] Each such inclusion concerns the smaller
lattice in the same generic lift space. A bound for
\(\sum_p J(j,p)\log p\) transfers to the weighted coefficient
places by the local degree identity
\[
 \sum_{w\mid p} e_w\,\frac{\log\mathrm Nw}{[K:\Q]}=\log p .
\] It remains to sum these lift exponents. Their contribution comes from lattice inclusion and is accounted for separately from the cohomology numerator.
\label{cor:division-length-6}
\label{cor:division-length-7}

\end{proof}

\paragraph{Aggregation of the least exponents.}
For the selected lift \(j\) let \(\mathcal C_w,\mathcal D_w\)
be the two local section lattices defining \(J_{j,w}\).
Both use the selected level, good subspace and tuple
\(\mu(w)\). The weight record of \(j\) uses the actual
weight at each embedding, up to an additive shift.
Let \(B_j(p)\) be the nonnegative integer obtained by the local
exponent rule below. The inclusion to be proved is
\[
                     p^{B_j(p)}\mathcal C_w\subseteq\mathcal D_w .
\]
Outside the active range the comparison uses the sum
of the full Weyl exponents over the \(f\) labels, with
the tail formula at the upper end. Inside the active
range the simultaneous local-model inclusions imply this
inclusion at the same rational-prime power. In those
models the indices with larger weights use the smaller of the full Weyl exponent
and the floor of the coarse bound for the surviving depths, and the indices with smaller weights
keep the elementary exponent.

\begin{lemma}[Summing the least local exponents]
In the specified size, jump and exceptional-prime range,
the least local exponents satisfy, uniformly in the lift and the coefficient field,
\[
 \sum_{w\ {\rm twisted}}\omega_wJ_{j,w}\le32\mathcal E,\qquad\omega_w=\frac{\log\mathrm Nw}{[K:\mathbb Q]}.
\]
\end{lemma}
\begin{proof}
Using \(p\) as a scalar costs \(e_w\) per power, so
\(J_{j,w}\le e_w B_j(p)\). Put \(\omega_w=\log\mathrm Nw/[K:\Q]\)
and let \(P_{\rm tw}\) be the rational primes below the
twisted places. Then
\begin{align}
 \sum_{w\ {\rm twisted}}\omega_w J_{j,w}
 &\le\sum_{p\in P_{\rm tw}}B_j(p)
            \sum_{\substack{w\ {\rm twisted}\\ w\mid p}}e_w\omega_w
 \le\sum_{p\in P_{\rm tw}}B_j(p)\log p \notag\\
 &\le 30|T_2|mbz^2+2mb
 \le32\mathcal E . \label{eq:division-least-sum}
\end{align} For the third inequality, use the size, jump range and exceptional-set conditions of the weight record in the prime-sum bound. Its
main and exceptional contributions together are at most
\(30|T_2|mbz^2\), and its tail contributes at most \(2mb\).
Here \(T_2\) bounds the number of indices with larger weights and has
size at most \(f\).

The last step uses \(z\ge1\) and
\(B_0=fmbz\). The estimate has the same constant
for every lift and is independent of the
coefficient degree. The calculation just given transports the prime sum.
Sections~\ref{subsec:lift-exponent-models} and
\ref{subsec:lift-exponent-primes} give the cover, weight
comparison and prime sum for this record, with the
representation-theoretic hypotheses used for the local linkage
and the integral graded bound. The paired row-index remainder
enters separately in Section~\ref{sec:word-division}.
\label{cor:division-length-8}

\end{proof}
\subsubsection{Lattice inclusions on the selected lift}\label{subsec:lift-exponent-models}
\hypertarget{source-lift-exponent-models}{}

The least exponents entering the signed index measure a further lattice inclusion beside the congruence length. To bound them, put the two coefficient lattices in the same space of forms at the selected lift.

Fix a selected lift, a tuple and a coefficient place \(w\mid p\).
Write \(Z\) for the selected space of forms at its row level.
In the coordinates of \(Z\), \(\mathcal C\) and \(\mathcal D\)
are the inverse images of the value lattices for the dual Weyl
and Weyl coefficients, respectively. The coordinate maps on
the finite class set identify both lattices simultaneously.
\begin{proposition}[Weyl inclusion on the original section lattices]
\label{prop:local-weyl-value-inclusion}
On the selected tuple and place \(w\mid p\), the sum of the exact Weyl exponents at all labels gives
\[
 p^{\sum_\tau e_p(\lambda_\tau)}\mathcal C_w\subseteq\mathcal D_w.
\]
For \(y<p\le h\) as in Lemma~\ref{hyp:weyl-prime-range},
exclude \(q,d,p_{\rm de},p'_{\rm de}\) and primes dividing an
integer of \(\mathcal I\). A power that gives the value inclusion at every local model under
an occurring integral Hecke character gives the same inclusion on
\(\mathcal C_w,\mathcal D_w\) at the original level.
\end{proposition}

\begin{proof}
\proofstep{Tensoring the exact inclusions and localizing the value map.}
At each coefficient label \(\tau\) choose a local basis of the
underlying lattice. The two coefficient models in that basis
use the dual Weyl and Weyl lattices \(W_\tau^\vee,W_\tau\) in
the prescribed rational realizations. The inclusion of
Section~\ref{subsec:lift-exponent-primes} tensors to
\[
 p^{\sum_\tau e_p(\lambda_\tau)}
          \bigotimes_\tau W_\tau^\vee
       \ \subseteq\  \bigotimes_\tau W_\tau .
\]
An additive weight translation twists both sides by the same
determinant line and does not change this exponent. If
\(\operatorname{ev}\) is the finite map of value coordinates,
the inverse image of the Weyl lattice is the kernel of its
composition with reduction modulo that lattice.
Writing \(q\) for reduction modulo the Weyl value lattice, exactness of localization gives
\[
 (\ker(q\circ\operatorname{ev}))_w=\ker(q_w\circ\operatorname{ev}_w).
\]
Equivalently one can
multiply the finitely many denominators outside \(w\) occurring
in the individual coordinates. The displayed inclusion therefore
gives the asserted exponent on each tuple.
\label{cor:lift-exponent-models-1}

\proofstep{Comparing the two lattices in the same operator space.}

We describe the sharper bound at primes in the weight range
\(y<p\le h\) outside the exceptional interval and specified
set, using the tuple assigned to \(w\). The base level and selected block hypotheses give finite free tensor value lattices over the same local coefficient ring. Choose
an embedding of the coefficient field into \(\overline{\Q}_p\)
inducing \(w\), and adjoin the roots of unity for the finite
level character to the induced local field. Its integer ring
\(O\) is a DVR. The integral frame data are changes of basis
on the two local value lattices, including the level group
matrices and character.

In the common interpolation frame a Hecke correspondence is represented on both lattices by the same operator
\[
 T=\sum_j a_j R(g_j),\qquad a_j\in O,
\]
with \(R(g_j)\) the specified right translations. Thus on the auxiliary space \(M\) of
integral forms the operators define an action of the commutative
polynomial algebra \(\mathcal T\) over \(O\) in these operators.
This action preserves the forms with values in the Weyl lattice,
whose submodule is \(M_{\rm W}\).

\proofstep{Detecting a failed inclusion at an occurring residual character.}
Take the integral form \(x_0\) from the dual Weyl reading of a
vector of \(\mathcal C\). Its value coordinates in the auxiliary
model satisfy covariance: it holds after scalar extension, and
extension is injective on the coordinates. The cyclic module
\[
 C=\mathcal T x_0\subset M
\]
is finite free over \(O\), and its
generic extension lies in the selected Hecke-stable span.
Write \(z=[p^a x_0]\) in \(M/M_{\rm W}\). If \(z\ne0\),
choose a maximal ideal \(\mathfrak m\) containing its annihilator.
Then \((\mathcal T z)_{\mathfrak m}\ne0\).
The kernel of \(\mathcal T\to\operatorname{End}_O C\)
kills \(x_0\) and hence \(z\). Thus the inclusions of ideals
\[
     \ker(\mathcal T\to\operatorname{End}_O C)
                   \subseteq\operatorname{Ann}_{\mathcal T}(z)
                   \subseteq\mathfrak m
\]
identify \(\mathfrak m\) with a maximal ideal of the
finite image algebra \(A\).

In a splitting field simultaneously triangularize the action on
\(C\). The diagonal characters \(\chi_j\) are integral,
and the matrices of \(J=\bigcap_j\ker\chi_j\) are strictly upper
triangular. In a triangular basis, with \(C_i\) the span of the first
\(i\) vectors, this gives
\[
 JC_i\subseteq C_{i-1}\quad(1\le i\le\dim C),
 \qquad J^{\dim C}C=0 .
\]
The action of the image algebra is faithful, so \(J\) is nilpotent. The product of integer rings
of the \(\chi_j\) is integral over \(A/J\). Lying over
gives a coordinate character with reduction \(\mathfrak m\).
Its generalized eigenspace is nonzero. On the generalized eigenspace \(V\), let \(a\) be the last exponent with \(\mathfrak n^aV\ne0\), for the maximal ideal \(\mathfrak n\) of the character. Then
\[
 0\ne\mathfrak n^aV\subseteq\ker\mathfrak n,\qquad \mathfrak n^{a+1}V=0.
\]
Any nonzero vector in this term is an eigenvector. It is an eigencharacter in the selected span:
extension of the interpolation frame preserves independence
because its nonzero full-rank minors remain nonzero.

\proofstep{Descent of the Weyl value test to the given row.}
Use the local models at all maximal ideals arising in this
way, equivalently those under integral eigencharacters occurring
in the selected span. If \(p^a\) gives the Weyl inclusion in
every such model, a maximal ideal detecting a nonzero class of
\(p^a x_0\) would contradict that inclusion. Thus \(p^a x_0\) has
Weyl values. Descent of the value tests gives the inclusion on
\(\mathcal C,\mathcal D\). In a localized coefficient basis the
condition is precisely nonnegative valuation on each coordinate.
It follows at the original level by the inverse-image description
above. This argument requires the same value and operator readings on both lattices.
\label{cor:lift-exponent-models-2}

\end{proof}

\paragraph{The two weights seen at a maximal ideal.}
At the chosen label put \(r=k_c+j+t+q\) and number the
coordinates from zero. Let \(A\) be the coefficient algebra
generated by the level matrices and the individual row-weight
projectors in the tensor. For a simple factor \(L\) of the
reduction, take decreasing real numbers \(\xi_{\tau,a}\) at
each label \(\tau\) and choose a joint weight \(\mu\) maximizing
\(\sum_{\tau,a}\xi_{\tau,a}\mu_{\tau,a}\) on \(L\).

If
\(\mu_{\tau,a}<\mu_{\tau,b}\) for \(a<b\), the transposition
of rows changes this sum by
\[
        (\xi_{\tau,a}-\xi_{\tau,b})
                     (\mu_{\tau,b}-\mu_{\tau,a})>0 .
\]
The row permutation preserves the weight set, so the
maximizer is dominant. The highest-weight theorem for this
coefficient algebra embeds \(L\) into the reduced dual Weyl
tensor of weight \(\mu\) and preserves the projectors. At one
label the theorem and the weight bounds give
\[
 \sum_a\mu_a=\sum_a\lambda_a,\qquad
 \lambda_{r-1}\le\mu_a\le\lambda_0,\qquad
 U_a=\lambda_a+r-1-a,\quad U_a'=\mu_a+r-1-a .
\]
In particular \(U,U'\) are strictly decreasing, of equal sum,
with entries in \([0,h_0]\), using \(\lambda_0+r-1\le h_0\).
\label{cor:lift-crystalline-depth-1}

Use the central Capelli generating polynomial
\(C_\nu(X)=\prod_{a=0}^{r-1}(X+\nu_a-a)\).
Its polarizations are coefficients of elementary row additions
followed by projection to the changed weights. The common degree
bound \(|\lambda|+|\mu|\) lets one use the same coefficient
operators on the reduced source and image. On a nonzero image
of \(L\) the two Capelli actions therefore have the same scalar
at every coefficient. The scalars are integral, so this gives
\[
       C_\lambda(X+r-1)\equiv C_\mu(X+r-1)\pmod p,
       \qquad \prod_a(X+U_a)=\prod_a(X+U_a')\ \text{in }k[X].
\]
Use the residue field \(k\) of the coefficient DVR here. Equality
of the two products gives the same multiset of residues, including
when several labels have the same residue.
\label{cor:lift-crystalline-depth-2}

Suppose the forms with values in \(L\) survive at the residual
ideal \(\mathfrak m\). At a chosen double coset \(c\), lift a
reduced value to the integral \(\mu\)-lattice. Its stabilizer
\(H_c\) has order invertible in \(O\). Averaging with the
level action \(R_c\) gives the lift
\[
        v_c^{\rm inv}=|H_c|^{-1}\sum_{h\in H_c} R_c(h)v_c .
\]
The reduction of \(v_c^{\rm inv}\) is the required value, since
the reduced value satisfies the stabilizer equations. Extend these
lifts by covariance. Let \(M_\mu\) be the integral weight-\(\mu\) forms and \(\bar v\) the surviving residual eigenform. Commutation of Hecke translations with reduction gives
\[
 \operatorname{Ann}(M_\mu)\subseteq\operatorname{Ann}(\bar v)=\mathfrak m.
\]
The finite torsion-free Hecke algebra argument
above now gives an integral characteristic-zero eigensystem
over \(\mathfrak m\), realized by a weight-\(\mu\) form since the
coordinate maps preserve independence. A residue-field embedding
aligning the two residual systems lifts to a local Galois
automorphism. This gives the congruence on every generator.
\label{cor:lift-crystalline-depth-3}

\paragraph{Automorphic and crystalline realizations at the chosen place.}
Use the coefficient field and the twisted tuple of this
lift at \(w\mid p\). Choose a \(p\)-adic embedding \(j_w\) inducing \(w\) and an isomorphism \(\iota:\overline{\Q}_p\to\C\). On the coefficient field their composite is prescribed:
\[
 \iota\circ j_w=j_{\mathbb C},
\]
where \(j_{\mathbb C}\) is the selected complex embedding. The latter comparison is algebraic.
The finite Hecke operators can be read through it because
their matrices on the finite class model are formed by
the same right translations and transported level
matrices.

At each member of the selected CM half, use
the finite adelic reading extending that \(p\)-adic
embedding. In the adapted basis the level matrices and
their inverses have integral entries. The twist is a
finite character, so adjoining the roots of unity of its
order puts these matrices and character values in the
same ring of integers \(O\). Evaluation at the finite
double classes gives the common integral form model.
The selected tame level makes the stabilizer orders units
in \(O\), including for the smaller weight \(\mu\).
\label{cor:lift-frame-galois-1}

\begin{proposition}[The oriented crystalline frame and the marked string]
\label{prop:lift-frame-crystalline-reading}
In the selected definite model, with the prescribed algebraic
type and spherical local factors over \(p\), the attached
crystalline representation has the rank-\(r\) normalized
Satake readings. For the \(\lambda\) system assume the specified
length-\(t\) divisor of the Hecke polynomial, \(r<2t\), and the
isolated length-\(t\) run in the weight at every label. Its
unique long Speh factor gives the marked consecutive crystalline
string. After the common character twist the \(\mu\)-frame has the following jumps at the two labels, in the given orientation:
\[
 \text{jump at }\tau=h_0-U'_{\tau,\bar a},\qquad\text{jump at }c\tau=U'_{\tau,a}.
\]
The marked string uses this orientation and the Frobenius transport of the proof.
\end{proposition}

\begin{proof}
\proofstep{The Hecke action after transport to the definite group.}
The characteristic-zero coefficient model first gives
a vector-valued form on the original definite space.
The isometry to the standard definite group transports
its rational and level actions. Use the algebraic
coefficient action at infinity to extend the finite
class values. Pair them using the positive invariant
Hermitian form on the coefficient type and take a
nonzero irreducible component of the resulting weight
space, using an invariant orthogonal complement.
The transported finite Hecke sums are the same on these
components.

In the split reading the prime is the
conjugate of the marked row prime. An isometry of the
two split groups identifies the right coset sums using
the dual double-coset label, which sends the list of elementary-divisor exponents to
\[
 (a_1,\ldots,a_r)\longmapsto(-a_r,\ldots,-a_1).
\]
This reads the given
Hecke eigenvalue, under \(\iota\), at the conjugate
prime. At primes of \(E\) above \(p\) the selected
level is spherical and the level character is trivial
on the corresponding hyperspecial group. The operators
indexed by good split primes away from \(p\) include
all coefficients of their Satake polynomials.
\label{cor:lift-frame-galois-2}

\proofstep{Attachment to the normalized Speh factors.}
Use unitary base change
with these local readings, the discrete spectrum
classification for \(\mathrm{GL}_r\), and Galois
attachment for the regular algebraic polarized cuspidal
constituents. Writing \(F=E^+\) for the totally real field of the definite group, the field hypotheses are
\[
 [F:\Q]\ge2,\qquad E/F\text{ unramified at all finite places}.
\] Base change has the
archimedean infinitesimal characters of the selected
algebraic type and the spherical components over \(p\).
Its discrete factors are Speh representations with
conjugate self-dual unitary cuspidal data. The spherical
component of the induced local product forces the
cuspidal local factors to be spherical.

For a cuspidal datum \(\pi\) of rank \(a\) in a Speh factor
of length \(d\), put
\[
 s_j=\frac{r-1}{2}+\frac{d-1}{2}-j,
 \qquad c_j=\frac{a-1}{2}-s_j \qquad (0\le j<d).
\]
In the unitary infinitesimal parameter let \(z\) range over
the exponents of \(\pi\) at a label. Regular algebraicity of
the full rank-\(r\) parameter gives
\[
 s_j-z\in\Z,
\]
because
\(z+(d-1)/2-j\) occurs there. The same regularity makes the
cuspidal exponents distinct at every label. The twist
\(\pi_j=\pi\otimes |\det|^{c_j}\) has exponents \(z+c_j\)
in \((a-1)/2+\Z\), and is polarized by the norm-power
character descending from \(|\cdot|^{2c_j}\). If \(a\) is odd,
multiply the polarizing character on the real subfield by
the quadratic character of the CM extension to give the
required sign at every real place. Its pullback is trivial.

Apply
Barnet-Lamb--Gee--Geraghty--Taylor, Theorem~2.1.1, to this
regular algebraic polarized cuspidal representation
in rank \(a\). With unitary Satake roots \(\alpha_{\pi,h}\) and Hodge labels
\(z_{\pi,\tau,h}\), the reciprocal normalization in rank \(a\)
reads
\[
 \operatorname{HT}_\tau(\rho_{\pi_j})
     =\{(a-1)/2-z_{\pi,\tau,h}-c_j\}_h
     =\{s_j-z_{\pi,\tau,h}\}_h,\qquad
 \operatorname{Sat}(\rho_{\pi_j})=\{q^{s_j}\alpha_{\pi,h}\}_h
\]
at the indicated good split prime for the Satake expression.
The first equality is the Hodge--Tate normalization and the
second multiset uses \((a-1)/2-c_j=s_j\).
The representation is crystalline above \(p\) because
\(\pi_j\) is spherical there. At a good split prime of
norm \(q\), the twist and the rank-\(a\) local Langlands
normalization multiply the unitary Satake roots by
\[
 q^{-c_j+(a-1)/2}=q^{s_j}.
\]

The direct sum over the cuspidal data and \(j\) therefore
has the full rank-\(r\) Satake roots scaled by
\(q^{(r-1)/2}\), with the corresponding full set of
crystalline jumps. Purity in Theorem~2.1.1 gives equal
absolute values for the unramified roots in each
cuspidal factor. If the cuspidal roots in unitary normalization have the common absolute value \(A>0\), then their product has modulus
\[
 A^a=1,
\]
where \(a\) is the rank of this cuspidal factor. This equality uses the unitary central character. Thus \(A=1\).
\cite{labesse-cm,blgg-weight}\label{cor:lift-frame-galois-3}

\proofstep{Identifying the unique long rank-one string.}
For the selected system of weight \(\lambda\) there is
additional information. Its occurrence in the chosen
block of the row supplies, at a good split prime, the
divisor of its Hecke polynomial coming from the
length-\(t\) string. Its center has modulus one in
unitary Satake normalization. To identify the global
factor, each cuspidal Satake multiset in the discrete
decomposition has unit modulus at this prime, by
temperedness of the cohomological conjugate self-dual
cuspidal factors in their unitary normalization. If a Speh constituent has length \(s\) and cuspidal rank \(a\),
its largest root modulus is \(q^{(s-1)/2}\), and its rank
contribution is \(as\). A root of modulus \(q^{(t-1)/2}\)
therefore requires
\[
       s\ge t,\qquad as\le r<2t .
\]
Hence \(a=1\), and two such constituents cannot fit in rank \(r\).

At each
archimedean label the specified length-\(t\) constant
run in the highest weight has strict drops at its two
ends when those ends are present. The shifted
infinitesimal exponents at this label are distinct.
The long Speh constituent gives a consecutive interval
of exponents. It cannot cross either drop, because each
drop leaves an absent intermediate exponent. The other intervals together have at most \(r-s<t\)
entries. They cannot cover the designated consecutive run of
\(t\) entries. The long interval must meet that run and cannot
cross a missing exponent at either boundary, so \(s\le t\).
Combining the size bound for the constituent and the two boundary gaps gives
\[
 t\le s\le t,\qquad s=t.
\]
Thus the long interval is exactly the specified run. Its Galois factors are the \(t\) successive
crystalline twists of a rank-one character. The
remaining Speh factors give \(\rho\).
\label{cor:lift-frame-galois-4}

\proofstep{Filtration jumps for the two CM orientations.}
The common normalization fixes their readings on both
orientations. Write \(\tau\) for a member of the
\(p\)-adic half chosen by the frame, \(c\tau\) for its
conjugate, \(s_\tau\) for the determinant weight shift
there, and
\[
 \bar a=r-1-a
\]
for a reversed coordinate.
Before the common twist, the type on \(\tau\) has highest
weight \(\mu_{\tau,a}+s_\tau\) and the type on \(c\tau\)
has its dual. The algebraic Hecke character with
crystalline jumps \(h_0-r+1+s_\tau\) on \(\tau\) and
\(-s_\tau\) on \(c\tau\) is chosen unramified over \(p\).
The sum of the two jumps is constant on all pairs, as
required for the algebraic character. Under the
reciprocal Frobenius convention the untwisted jump is
\(\bar a-\lambda_{{\rm aut},\bar a}\) in either local
highest weight \(\lambda_{\rm aut}\).

Adding the character
jump gives the two arrays
\[
 J_{\tau,a}=h_0-U'_{\tau,\bar a},\qquad
 J_{c\tau,a}=U'_{\tau,a},\qquad
 U'_{\tau,a}=\mu_{\tau,a}+r-1-a .
\]
Thus the marked orientation for \(U'\) is the conjugate
of the half used for the complex automorphic frame.
Apply the same twist to the selected \(\lambda\) system.
On that orientation the marked line jumps are
\(B_\tau+s\), \(0\le s<t\), with \(B_\tau\) the bottom
of the shifted run of \(\lambda\). On the other orientation
their bottom is \(h_0-B_\tau-(t-1)\), with the order of
the marked indices reversed. Both arrays and both
families of lines therefore lie in \([0,h_0]\).
The Frobenius permutation acts on all embeddings of the
unramified host and commutes with their CM involution.
This determines the transported label and its
orientation together in the crystalline comparison.
\label{cor:lift-frame-galois-5}
\end{proof}


\begin{lemma}[Identifying the marked factors in the companion reduction]
\label{lem:companion-residual-factors}
On the common coefficient extension of the two normalized
Hecke systems, the residual semisimplification is
\[
 \bar r_\mu^{\rm ss}\simeq \bar\rho^{\rm ss}\oplus\bigoplus_{s=0}^{t-1}\bar\chi_s.
\]
One may keep the marked occurrence of each scalar line in a
composition series of its actual reduced lattice.
\end{lemma}

\begin{proof}
At the good degree-one Frobenius classes, congruence of
the two Hecke systems gives
\[
 \det(X-r_\mu(g))\equiv
 \det(X-\rho(g))\prod_{s=0}^{t-1}(X-\chi_s(g))
                 \pmod {\mathfrak m_{\overline{\Q}_p}} .
\]
Indeed the reversed polynomial has Hecke coefficients times
signed integral norm powers and powers of the common twisting
character, which has unit values by compactness. For degree \(r\), reversing by
\[
 P(X)\mapsto X^rP(X^{-1})
\]
is a permutation of the coefficient list. It therefore
preserves the congruence after the common unit normalizations. The congruence
locus is closed because all coefficients are continuous and
the open unit ball in an ultrametric field is closed.
Degree-one Chebotarev density after finitely many omissions
extends it to every group element, also for the reciprocal
Frobenius convention. Choose a common finite coefficient field
and stable lattices.

For the two residual representations \(\bar R_1,\bar R_2\), Brauer--Nesbitt gives
\[
 \det(1-X\bar R_1(g))=\det(1-X\bar R_2(g))\ (\forall g)\ \Longrightarrow\ \bar R_1^{\rm ss}\simeq\bar R_2^{\rm ss}.
\]
A composition series therefore has the
prescribed line characters among its factors with their
multiplicities. Keep their positions in the filtration.
The block comparison takes these lattices over a common
finite coefficient extension with the marked characters
defined there. Compactness gives an invariant lattice
by taking the span of translates of one integral basis.
The coefficient field can be enlarged to contain the
labeled fields and the primitive \(p\)-th roots used in
the cyclotomic frame. This computation uses the original
residual filtration on the \(r_\mu\) side, with its
marked factors identified by the semisimplification.
\label{cor:lift-crystalline-depth-4}
\end{proof}

\paragraph{What the residual lines imply about the head.}
\hypertarget{source-lift-weight-exchange}{}
Let \(H=h_0\) and \(K=\lfloor H/(p-1)\rfloor\).
Here \(p>y>8\), \(p\) is unramified in the CM field, and
the large-prime inequalities give \(t+2K<p\).
Use all embeddings of the CM field at \(p\). At the conjugate embedding use the decreasing sequence obtained by reflection and reversal. In zero-based coordinates the operation is
\[
 (u^c)_a=H-u_{r-1-a}\quad(0\le a<r),\qquad (u^c)^c=u .
\] The
bottom of the new block is \(H-(B+t-1)\). Write
\(\sigma i\) for the source label of the Frobenius matrix
ending at \(i\).

In the local cyclic indexing it is \(i-1\).
This permutation commutes with conjugation. The Galois attachment can use either
orientation for the whole family. Pulling back the two
sequences and the block along that involution identifies its
\(J\) with \(U'\), and its line weights with
\(b_{i,s}=B_i+s\) at each embedding \(i\). A calculation at
the transported label then reads at \(i,\sigma i\) in the
original labels. Reflection preserves the characteristic-polynomial comparison. Indeed, for any list of residues of length \(r\),
\[
 \prod_{a=0}^{r-1}\bigl(X-(H-u_a)\bigr)=(-1)^r\prod_{a=0}^{r-1}\bigl((H-X)-u_a\bigr).
\]

\begin{proposition}[A depth bound from the residual lines]\label{prop:residual-head-depth}
For the shifted highest weights, conjugate labels and consecutive
line block just specified, the deficit at a head cut satisfies
\[
 \sum_{a<k_c+j}(\lambda_{i,a}-\mu_{i,a})
   \le qH+(p m_{\sigma i}+m_i)(t+2K).
\]
Here \(m_i\) counts the original shifted weights outside the consecutive
block at \(i\) whose residues occur in the block enlarged by \(K\)
at each end. The count \(m_{\sigma i}\) uses the Frobenius source label. At the indices satisfying the shape and
prime exclusions of Section~\ref{subsec:lift-exponent-primes},
both corrections vanish and the depth is at most \(\lfloor qh\rfloor\),
uniformly over the occurring Hecke ideal and the weight \(\mu\).
\end{proposition}

\begin{proof}
Use the norm coordinate of Section~\ref{subsec:cut-wach-models},
written \(X=Z\) in this calculation, with the same terminal
Hodge labels. Apply its construction to
the lattice of the congruent weight \(\mu\), to the inverse
images of the residual flag and to the specified scalar lines.
Their coefficient extension is over the selected place. Their
jumps are respectively \(U'_{i,a}\) and \(b_{i,s}\) and lie
in \([0,H]\). Write \(F\) for the residue field. In the resulting
\(F[[X]]\) lattice the reduced Frobenius has an integral right
inverse up to \(X^H\). For any \(d\)-rowed minor \(M\) of the reduced Frobenius matrix the elementary divisors give
\[
 \operatorname{ord}_X M\ \ge\ s_i(d),\qquad s_i(d)=\sum_{a=r-d}^{r-1}U'_{i,a}.
\]
The first cyclotomic matrix coefficient \(\Theta_i\) has
characteristic roots \(-\bar c U'_{i,a}\).
The residual filtration gives the same ranks and marked lines in
\(F((X))\) by exact 'etale realization.

Its intersection with
\(F[[X]]^r\) is a saturated flag, so extending bases at successive
steps gives one basis adapted to it. Each marked line occupies its own
diagonal block. Its two chosen generators in the Laurent
line differ by a unit times \(X^{\kappa_{i,s}}\), with the
sign fixed so that Frobenius in the adapted basis has order
\[
        \alpha_{i,s}=b_{i,s}+p\kappa_{\sigma i,s}-\kappa_{i,s}.
\]
In fact Frobenius sends the parameter to \(X^p\), and the
reduced rank-one Frobenius is a unit times \(X^{b_{i,s}}\).
The latter follows from the rank-one determinant, whose
order before reduction is \(b_{i,s}\) at \(X-p\) by height
and the local elementary divisor. Integrality gives
\(\alpha_{i,s}\ge0\) and the height inverse gives
\(\alpha_{i,s}\le H\) on the singleton block. The inverse
is block triangular after inverting \(X\), so the two
diagonal entries multiply to \(X^H\).

\proofstep{From the determinant to the minor estimate.}
To constrain individual head weights we pass from the determinant to a minor.

For the cyclotomic generator, \(\gamma(X)/X\) has first
terms \(1+\bar c X\) modulo \(X^2\) in this reduced
coordinate. For a unit \(v=v_0+v_1X+O(X^2)\), the constant and linear terms cancel in the quotient:
\[
 \gamma(v)\equiv v_0+v_1X\equiv v\pmod{X^2},\qquad \gamma(v)/v\equiv1\pmod{X^2}.
\] The coefficient of \(X\) on the same marked line
is therefore \(-\bar c(b_{i,s}-\kappa_{i,s})\).
An integral basis change conjugates this first coefficient by
its constant matrix, since the cyclotomic action fixes the
change of basis modulo \(X^2\).

Taking characteristic
polynomials of the block triangular matrix and cancelling
the nonzero common scalar shows that the residues
\(b_{i,s}-\kappa_{i,s}\), counted with multiplicity, occur
among those of \(U'_i\). The determinant of the Frobenius
minor on the marked positions is the product of their
diagonal entries. The minor divisibility then gives
\begin{equation}\label{eq:lift-line-minor}
       s_i(t)\le\sum_s b_{i,s}
              +p\sum_s\kappa_{\sigma i,s}-\sum_s\kappa_{i,s}.
\end{equation}
The flag and the rank-one comparison here come from the
same residual representation at the same maximal ideal.
\label{cor:lift-weight-exchange-1}

Since \(b_{i,s}\) and \(\alpha_{i,s}\) belong to \([0,H]\),
we have \(|p\kappa_{\sigma i,s}-\kappa_{i,s}|\le H\).
For fixed \(s\), let \(j\) maximize \(\lvert\kappa_{j,s}\rvert\).
Apply the relation at \(i=\sigma^{-1}j\), where the term
with coefficient \(p\) is \(\kappa_{j,s}\). The triangle inequality at this predecessor gives
\[
 p|\kappa_{j,s}|\le H+|\kappa_{\sigma^{-1}j,s}|
            \le H+|\kappa_{j,s}| .
\]
Hence every \(|\kappa_{i,s}|\le H/(p-1)\).
To bound the sum of the shifts, enlarge the block
\(I=[B_i,B_i+t-1]\) by \(K\) at both ends to \(I^+\).

Let \(m_i\) count original weights outside \(I\) whose
residue occurs in \(I^+\). The strict length bound makes
reduction modulo \(p\) injective on the integers in \(I^+\).
The \(t\) shifted entries \(b_{i,s}-\kappa_{i,s}\) lie
there. Match their residues to distinct occurrences among
the \(U_i\) using the two multiset comparisons. At most
\(m_i\) are matched outside \(I\). All other matches are
equalities in \(I^+\). Pair the remaining shifted entries
arbitrarily with the unused block entries. Each difference
is bounded by \(t+2K\), so
\[
                \left|\sum_s\kappa_{i,s}\right|
                            \le m_i(t+2K).
\]
This also holds at \(\sigma i\), where the exceptional count
uses its own block. Reflection preserves the count because it
carries the block, its enlargement and residue congruence
bijectively to their reflected versions.

\proofstep{Returning to the head of the weight.}
The deficit of the first \(k_c+j\) weights equals the excess
of the rest by equality of total sums. Of the last \(t+q\)
entries of \(U'_i\), the smallest \(t\) have sum \(s_i(t)\)
and the other \(q\) are at most \(H\) apiece. In \(U_i\)
the corresponding total is at least \(\sum_s b_{i,s}\)
because the bottom \(q\) entries are nonnegative. Equation~\eqref{eq:lift-line-minor}
thus proves
\[
 \sum_{a<k_c+j}(\lambda_{i,a}-\mu_{i,a})
   \le qH+(p m_{\sigma i}+m_i)(t+2K).
\]
The staircase cancels term by term on the left. The shape
and prime exclusions of Section~\ref{subsec:lift-exponent-primes}
give \(m_i=m_{\sigma i}=0\). Every surviving highest
weight at this index with a larger weight therefore has depth at most
\(i_i=\lfloor qh\rfloor\), with the choice of
\(\mu\) and of the occurring Hecke ideal unrestricted.
\label{cor:lift-weight-exchange-2}
\end{proof}


\paragraph{The localized tensor inclusion.}
For each embedding \(\tau\), let \(V_\tau=\nabla(\lambda_\tau)\)
be the dual Weyl module of highest weight \(\lambda_\tau\) over \(O\).
Write \(B_\tau=\Delta(\lambda_\tau)\subset V_\tau\) for its Weyl lattice.
Write \(V=\bigotimes_\tau V_\tau\).

\paragraph{Divisibility on one row-weight space.}
\hypertarget{source-lift-pieri-cost}{}
We give the bounded-depth computation over \(\Z_{(p)}\).
For a dominant weight \(\nu\), a row weight \(v\) and a vector
of coefficients \(\xi=(\xi_a)_a\), put
\(d_\xi(v)=\sum_a\xi_a(\nu_a-v_a)\).
In the application, \(\xi_a=1\) on the first \(k_c+j\) coordinates
and \(\xi_a=0\) on the others.

\begin{lemma}[Rectangular Pieri denominator]
\label{lem:rectangular-pieri-denominator}
Consider first a rectangular
initial block \(\nu=(w^b,\beta)\), where \(\beta\) is a
partition of length at most \(d\) and \(\beta_1\le w\).
An upper bound for the extra exponent in passing from the
last \(d\) coordinates to the full coefficient module at
\(n=bw-\sum_{a<b}v_a\) is
\begin{equation}\label{eq:lift-pieri-polynomial}
 v_p(D_n(w)),\qquad
 D_n(Z)=\prod_{c}(Z-c)^{m(c)} .
\end{equation}
Here one takes partitions \(\tau\supset\beta\) with
\(|\tau/\beta|=n\), column lengths of \(\tau/\beta\) at
most \(b\), and no added box of content \(w\). Content is
column minus row, both numbered from zero. The integer
\(m(c)\) is the maximum number of boxes of content \(c\)
in such a skew diagram, with empty maximum zero.

In
particular \(D_n(w)\ne0\). If \(p^e\) suffices on the
Levi weight spaces with \(d_\xi\le I\) (extend their weights
by the initial rectangle), then \(p^{e+v_p(D_n(w))}\)
suffices on the indicated full weight space whenever
\(d_\xi(v)\le I\) and \(\xi\) is decreasing in the weak sense.
\label{cor:lift-pieri-cost-1}

\end{lemma}

\begin{proof}
\proofstep{Integral raising and lowering tests.}
Use the polynomial representation over \(\Q\) with the
contravariant Fock form normalized to give the highest
weight vector norm one. The Levi highest subspace \(W\)
is generated from this vector by the last \(d\) rows. The
first block acts there by the determinant character and the
upper cross-block operators kill it. Let \(F_{u\alpha}\)
be the cross-block lowering operator from first row \(\alpha\)
to last row \(u\), and \(E_{\alpha u}\) its adjoint.
The lowering operators commute with each other. Their products
of order \(n\) applied to \(W\) span the depth \(n\) subspace.
Indeed reordering words in row operators puts the lowering
operators on the left and the raising operators on the right.
Let \(T\) be the sum of the first diagonal operators. On the span of depth \(n\) it acts by \(bw-n\). Since it is self-adjoint in the contravariant form, vectors at different depths satisfy
\[
 (n'-n)\langle v_n,v_{n'}\rangle=\langle Tv_n,v_{n'}\rangle-\langle v_n,Tv_{n'}\rangle=0 .
\]

For a multiset \(m\) of \(n\) pairs let
\(m!=\prod_{u,\alpha}m_{u\alpha}!\).
The lowering generators of the integral lattice are
\(F^m/m!\) on the Levi Weyl lattice. Each is a product of
integral divided row powers. By adjunction, the test against a divided lowering word is a test in the dual Levi lattice. For a Levi vector \(z\),
\[
 \left\langle v_0,\frac{F^m}{m!}z\right\rangle =\langle y_m,z\rangle,\qquad y_m=\frac{E^m}{m!}v_0.
\]
These entries determine \(v_0\), by the spanning assertion
and nondegeneracy. In the polynomial model, restriction to
the Levi block identifies its highest weight generator and
its divided row words with their smaller-rank counterparts.
The restricted dual tests are integral there because their pairings
with the transposed words read the coefficients on the unit
lower triangular chart. The ordered row additions extract
these coefficients and restriction is injective on \(W\).
\label{cor:lift-pieri-cost-2}

\proofstep{The commuting polynomial on multisets.}
To compute with the entries, represent a vector by an array
of ordered lowering words of length \(n\) with coefficients
in \(W\). Let \(L\) denote evaluation of the words and \(S\)
the map which sums the array over words of each multiset.
Thus \(L\) factors through \(S\). In tensor position \(j\),
numbered from one, set
\[
 J_j=\sum_{u,v}e^{(j)}_{uv}E^W_{vu}
                       +\sum_{a<j}(a\ j)_Y ,
 \qquad A(Z)=\prod_{j=1}^n(Z-J_j).
\]
Here \(e^{(j)}_{uv}\) acts on the last-block row index,
\(E^W_{vu}\) on the Levi coefficient and the transpositions
exchange just the last-block indices. The raising tests of
\(Lz\) satisfy
\[
                    \iota Lz=S A(w) z,\qquad
                    (\iota v_0)_m=E^m v_0/m! .
\]
For this identity commute the first raising operator past
each lowering in a word, using
\([E_{\alpha u},F_{v\beta}]
  =\delta_{uv}E_{\alpha\beta}-\delta_{\alpha\beta}E_{vu}\).

On the highest subspace the first term gives the scalar
\(w\), the second gives the Levi action. Commuting the
resulting row generators through the other letters gives
the transpositions in \(J_j\). Iteration over the letters
gives the ordered product on the sum over permutations of
the word. In the sum over all orders of a fixed multiset, the stabilizer permutes the repeated copies of each pair. Its cardinality is
\[
 |\operatorname{Stab}(m)|=\prod_{u,\alpha}m_{u\alpha}!=m!,
\]
giving the divided raising test above on collecting equal words.
These are identities over characteristic zero. The
coefficients of \(A(Z)\) descend to commuting operators on
multiset arrays. Indeed the successive partial quadratic
Casimirs commute, their differences give the \(J_j\), and
adjacent transpositions together generate the permutations.
Conjugating the defining sums by the adjacent transposition on the last-block indices gives the two relations
\[
 sJ_a=J_{a+1}s-1,\qquad sJ_{a+1}=J_as+1,\qquad s=(a\ a+1).
\]

The transposition \(s\) commutes with the other \(J_i\).
Using the two swap relations and commutativity of \(J_a,J_{a+1}\)
gives
\[
        s(J_a+J_{a+1})=(J_a+J_{a+1})s,\qquad
        sJ_aJ_{a+1}=J_aJ_{a+1}s .
\]
These are the two coefficients that show commutation with
\((Z-J_a)(Z-J_{a+1})\), hence with \(A(Z)\). Swaps on the
first-block indices commute with all the \(J_i\). Its coefficients preserve integral
arrays since the operator upstairs uses integer permutations
and Levi row operators and \(S\) is onto over the integers.

\proofstep{Contents on the surviving Pieri paths.}
Tensoring the irreducible Levi module successively with
the defining representation gives the rational Pieri paths
which add one box per step. On a path the eigenvalue of
\(J_j\) is the content of the box just added. In fact the
quadratic Casimir on shape \(\kappa\) has scalar
\(\sum_a \kappa_a(\kappa_a+d-1-2a)\). For the box in row \(a\) the new column has zero-based index \(\kappa_a\).
The change of the form, minus \(d\), is
\[
 (\kappa_a+1)(\kappa_a+d-2a)
       -\kappa_a(\kappa_a+d-1-2a)-d=2(\kappa_a-a).
\]
Here the content is \(\kappa_a-a\), equivalently the
column number \(\kappa_a+1\) in one-based column notation
minus one, then minus \(a\). The tensor Casimir identity gives twice \(J_j\)
for this operator difference.

To see which paths can survive
\(S\), fix the counts \(n_\alpha\) of the first-block letters.
The multiset arrays with these counts, as a module of the
last block in characteristic zero, are
\(W\otimes\bigotimes_{\alpha=1}^b
\operatorname{Sym}^{n_\alpha}(\Q^d)\).
Indeed their multiset basis is indexed by one monomial per
first-block letter. By the horizontal-strip Pieri rule, let \(\lambda^{(j)}/\lambda^{(j-1)}\) be the strip contributed by the \(j\)-th symmetric power.

In each column \(c\),
\[
 (\tau/\beta)'_c=\sum_{j=1}^b\bigl((\lambda^{(j)})'_c-(\lambda^{(j-1)})'_c\bigr)\le b,
\]
since each summand is zero or one. Here the left side counts the boxes of the skew diagram in that column. The column bound kills the image under \(S\) of every
\(\tau\) with a column of length \(>b\) in \(\tau/\beta\).
For a surviving path, each successive \(J_i\) reads the
content of its added box, so the scalar on its \(\tau\)-space is
\[
       P_\tau(Z)=\prod_{\square\in\tau/\beta}
                   (Z-\operatorname{cont}(\square)).
\]
The product depends only on the skew shape. Since \(S\)
is equivariant, \(S A(w)\) has image in the sum of the
surviving spaces with \(P_\tau(w)\ne0\).
\label{cor:lift-pieri-cost-3}

\proofstep{The common content polynomial clears the inverse.}
There is an integral way to invert this last operation up
to the scalar \(D_n(w)\). Divide \(D_n(Z)\) by the universal
monic polynomial \(Z^n+\sum_{a<n}c_a Z^a\) over
\(\Z[c_0,\ldots,c_{n-1}]\), and evaluate the quotient
at \(w\) with the commuting coefficient operators of \(A\)
in place of \(c_a\). Call the resulting integral operator
\(Q\).

On every path occurring in the image,
\(P_\tau\mid D_n\) by the content maxima. The specialized
remainder is zero. The polynomial-division identity on a surviving path thus specializes to
\[
 A(w)Q=D_n(w)\,\mathrm{id}
\]
on the image under consideration.
This division incurs no denominators for spectral
projectors. The coordinate of a raising test has weight
\(v+\sum m_{u\alpha}(e_\alpha-e_u)\).
Each \(J_j\) respects the matching between coordinate
indices and their required weights, since its matrix unit
and Levi action have opposite weight changes. The depth condition is preserved since \(\xi_\alpha\ge\xi_u\) at each contributing pair. The weight shift in a raising coordinate has pairing
\[
 \left\langle\xi,\sum_{u,\alpha}m_{u\alpha}(e_\alpha-e_u)\right\rangle=\sum_{u,\alpha}m_{u\alpha}(\xi_\alpha-\xi_u)\ge0 .
\]
Thus under \(Q\) each coordinate remains in the dual Levi lattice at the allowed depth.
Multiplication by \(p^e\) puts these coordinates in the
Levi Weyl lattice.

Apply the lowering-word evaluation
to the multiset array \(p^e Q\iota v_0\). It belongs to
the full Weyl lattice because each undivided word is an
integer multiple of its divided word. Its raising tests
are \(p^e D_n(w)\iota v_0\), so injectivity of the tests
proves \eqref{eq:lift-pieri-polynomial}. The scalar prime
to \(p\) is a unit. The descent to the smaller rank here
uses the indicated polynomial Weyl lattices with their
normalized pairing.\label{cor:lift-pieri-cost-4}

\end{proof}

\paragraph{Counting the factors at the selected weight.}
\begin{lemma}[Depth bounds at the selected rectangular block]
\label{lem:selected-pieri-depth}
On the selected row-depth cutoff \(I=i_i\), write \(c_{\rm out}\) for the cost of the two parts outside the rectangle. Then
\[
 c_{\rm out}\le\frac{(j+q)(h_0-1)}{p-1}.
\] The rectangular factors are counted by
the shifted intervals below, for each integer \(0\le n\le I\).
In the nonexceptional range \(h_0<p^2\) this gives the stated
coarse exponent for the projected inclusion, using the original
Frobenius label and the same grading.
\end{lemma}

\begin{proof}
\proofstep{Complementary coordinates and the one-row costs.}
For a rectangular block at the end reverse the coordinates
and take complementary weights about a common bound \(M\).
This is contragredience with a determinant twist. It
interchanges the two block positions, replacing \(\xi\) by
its negative reverse and preserving the graded deficit.
In a one-row step the added skew diagram is a horizontal
strip. Two boxes of the same content, farther right in a
lower row, would force a common added column. Thus all
contents in this strip are distinct. The surviving boxes cannot reach content \(w\)
since each row starts at content at most \(w\) and must
avoid \(w\).

Taking the maximum of the shifted widths, the surviving factors therefore satisfy
\[
 1\le w-c\le h_0-1.
\] More precisely
for an appended row \(a\), numbered from zero, the
complementary rectangle is \(M-\lambda_a\), and the
smallest possible content is \(1-a\) for \(a\ge1\).
With \(M=\lambda_0\), all factors are at most
\(\lambda_0+a-\lambda_a-1\le h_0-1\).
In each such step the distinct positive integers \(w-c\)
lie in \([1,h_0-1]\), so their product divides \((h_0-1)!\).
The empty Levi contributes the product one. It follows that
\[
  v_p\!\left(\prod_c(w-c)\right)
  \le v_p((h_0-1)!)
  \le\left\lfloor\frac{h_0-1}{p-1}\right\rfloor .
\]
\label{cor:lift-pieri-cost-5}

Start with the full exponent of the first \(k_c\) rows
(zero for the selected nonexceptional prime). Append the next
\(j\) rows one at a time, using the zero grading. Append
the rectangle of length \(t\) at once, with
\(w=\lambda_{k_c+j}\) before reversal. For this step the
block depth equals the deficit above the rectangle by
equality of the two total weights. Finally append the \(q\)
bottom rows individually, retaining the grading by the
first \(k_c+j\) entries. The preceding argument with
decreasing \(\xi\) keeps the same allowance \(I=i_i\)
through these last steps. The middle step must therefore be bounded at precisely the integer depths
\[
 n\in\{0,\ldots,I\},\qquad I=i_i .
\] The two outside
steps together cost at most \((j+q)(h_0-1)/(p-1)\).

\proofstep{Which labels can occur at bounded depth.}
On the reversed Levi let \(M=\lambda_0\) and
\(\beta_b=M-\lambda_{k_c+j-1-b}\), with rectangle height
\(M-w\). For a box of content \(c\), put
\(x=M+t-1-c+q\). Subtract the top weight \(B_\tau=w+t-1+q\) of the original shifted block. The reversed height then contributes the same factor:
\[
 x-B_\tau=(M+t-1-c+q)-(w+t-1+q)=(M-w)-c .
\] The labels in a row start
one below its shifted weight and decrease by one at each
box. They are strictly above \(B_\tau\) by the omitted
content condition.

The first \(k_c\) original rows have
shifted weights in \([x_-,x_+]\) and the next \(j\)
at most \(x_-\). The latter rows supply at most \(j\)
boxes of each label below \(x_-\), none above or at that
endpoint. The first \(k_c\) rows supply at most \(k_c\) boxes of a
label, and at most \(t\) since boxes on one content
diagonal force a common added column. For \(x<x_-\) each
of these rows reaching \(x\) has at least \(x_--x\) added
boxes. Their total is at most \(n\). If \(c_x\) counts these rows reaching \(x<x_-\), then
\[
       c_x(x_--x)\le n\le I,\qquad
       c_x\le \min\{k_c,t,I/(x_--x)\}.
\]
This multiplicity estimate applies at each such label.

\proofstep{Counting multiples of the residue prime.}
The \(j\)-row contribution is bounded by
\(jv_p((x_--B_\tau-1)!)\le j(h_0-1)/(p-1)\).
All relevant positive \(x-B_\tau\le h_0<p^2\) have
valuation at most one. In \([x_-,x_+]\) the number
of multiples of \(p\) is at most \((x_+-x_-)/p+I_p\),
with the indicator from Section~\ref{subsec:lift-exponent-primes}.
For \(x<x_-\), use the representative of the residue of \(x_--B_\tau\) in \(\{1,\ldots,p\}\).

The positive residue representative is
\[
 r_p=1+\bigl((x_--B_\tau-1)\bmod p\bigr).
\] The
possible distances \(x_--x\) are \(r_p+ap\), strictly
less than \(x_--B_\tau\). The term \(a=0\) gives
\(\min(k_c,t,I/r_p)\) if this strict inequality holds.
For \(a\ge1\) sum \(I/(ap)\) to get at most
\((I/p)(1+\log(h/p))\) when \(p\le h\).
Enlarging the quantities in this bound gives the claimed coarse estimate:
\[
 k_c\le k,\qquad I\le qh,\qquad \frac{h_0-1}{p-1}\le\frac{h+r}{p-1}+6 .
\] The cost is an
integer, so one may take the floor. Finally the projected
inclusion over \(O\) follows from the integral standard
bideterminant spanning theorem. Each generator is homogeneous
in the row degrees, and its projection is either itself or
zero. Its prime-to-\(p\) clearing integer from the integral
computation becomes a unit in \(O\).
\label{cor:lift-pieri-cost-6}
\end{proof}


\begin{lemma}[Descent from projected Weyl values]
\label{lem:localized-projected-weyl-values}
With the integral weights of the tensor \(\bigotimes_\tau
\nabla(\lambda_\tau)\), use the projected inclusions of
Lemmas~\ref{lem:rectangular-pieri-denominator}
and~\ref{lem:selected-pieri-depth}. The Weyl inclusion on the actual forms at
a selected residual ideal is obtained by clearing each factor
by its full exponent, or by the smaller allowable depth cost
at a larger-weight factor. Write \(d_\tau\) for the charge so chosen at each label, \(V\) for the unrestricted tensor lattice and \(D_\varnothing\) for the lattice using the Weyl factor at every label. After passing to the forms at the selected residual ideal \(\mathfrak m\), the inclusion is
\[
 p^{\sum_\tau d_\tau}\,\mathcal F(V)_{\mathfrak m}
 \subseteq \mathcal F(D_\varnothing)_{\mathfrak m},
\]
where \(\mathcal F(L)\) denotes the module of forms with values in \(L\).
\end{lemma}

\begin{proof}
For a set \(S\) of
unrestricted labels, let \(D_S\) be the image of the tensor
using \(V_\tau\) on \(S\) and \(B_\tau\) elsewhere. Each
\(D_S\) is stable under the coefficient algebra and its forms
are stable under the Hecke translations. The exact exponents
\(e_\tau\) give \(p^{\sum e_\tau}V\subset D_\varnothing\)
by multiplication on pure tensors. Thus \(V/D_\varnothing\)
has finite length. Every simple factor of a subquotient of
this quotient occurs in the reduction of \(V\) modulo the
uniformizer \(\varpi\). Indeed, multiplication by a power of \(\varpi\) gives the surjection onto each successive quotient:
\[
 V/\varpi V\ \twoheadrightarrow\ \varpi^a V/\varpi^{a+1}V,\qquad [v]\longmapsto[\varpi^a v]\quad(a\ge0).
\]

Write \(P_{\tau,w}\) for the integral weight projector,
\(i_\tau\) for the depth allowance and \(c_\tau\) for the
integral bounded-depth cost just obtained. For
\(a\notin S\), set \(D=D_S\) and \(D'=D_{S\cup\{a\}}\).
If the full exponent is used at \(a\), then
\(p^{e_a}D'\subset D\) on pure tensors. At an index with a larger weight
where \(c_a<e_a\), form the actual intermediate lattice
\[
                          E=D+p^{c_a}D'.
\]
For any weight \(w\) at depth at most \(i_a\) we have
\(p^{c_a}P_{a,w}V_a\subset B_a\). Applying that projector
in just factor \(a\) shows \(P_{a,w} E\subset D\).

In a simple factor \(L\) of \(E/D\), the projector at its
highest weight \(\mu_a\) has nonzero image. This follows
from the highest-weight line in the coefficient classification.
If the highest weight had depth at most \(i_a\), the integral inclusion would give
\[
 P_{a,\mu_a}E\subseteq D\quad\Longrightarrow\quad P_{a,\mu_a}(E/D)=0,
\]
contradicting the nonzero image on \(L\). Thus all simple factors here have
excess depth. By the preceding crystalline estimate their
localized forms vanish. This argument requires stability of
the intermediate lattices and the integral action of the
projector on the full tensor.

Take a composition series for \(E/D\). At each successive
quotient map a form to the form obtained by projecting its
values to the simple factor. This map is Hecke-equivariant.
Filter \(E\) by its series \(D=D_0\subset\cdots\subset D_b=E\), and denote the projection at the last step by \(\pi_b\). The vanishing for the simple factor reads
\[
 (\pi_b\phi_b)_{\mathfrak m}=0,\qquad \pi_b:D_b\to D_b/D_{b-1}.
\]
There is \(s_b\notin\mathfrak m\) killing that image before
localization, so \(\phi_{b-1}=s_b\phi_b\) takes values in
\(D_{b-1}\). Composing the successive denominator clearings gives
\[
 (s_1\cdots s_b)\phi_b\ \text{with values in }D_0=D .
\]
Multiplying these denominators down the series proves the
localized inclusion of the forms with values in \(E\) into
those with values in \(D\).

In particular \(p^{c_a}\) moves
the forms on \(D'\) into the forms on \(D\) after localization.
Inducting over the unrestricted labels gives
\[
 p^{\sum_{a\ {\rm thick}}\min(e_a,c_a)+
       \sum_{a\ {\rm other}}e_a}
            M_{\mathfrak m}\ \subset\ (M_{\rm W})_{\mathfrak m}.
\]
All denominators used are units in this localization.
The simple-factor argument works separately at each integral
joint character over the ideal. The surviving bound is the
same for every such character in the prime range under
consideration. This proves the passage from the projected
coefficient inequalities to the lattice of actual forms.
\label{cor:lift-graded-inclusion-1}
\end{proof}



\begin{lemma}[Least exponent on the original Weyl value pair]
\label{lem:weyl-least-exponent-local}
For the coefficient models of Proposition~\ref{prop:local-weyl-value-inclusion},
let \(B_j(p)\) be the bound from Section~\ref{subsec:lift-exponent-primes}
for the weight record of lift \(j\). For the ramification index \(e_w\) and the least local exponent,
\[
 J_{j,w}/e_w\le B_j(p)\qquad(w\mid p).
\]
This normalized upper bound descends through the coefficient extension.
\end{lemma}

\begin{proof}
The exponent in these models is computed with the actual
Frobenius and conjugation labels of the frame. In the indicated
range the two window correction counts vanish by the shape
calculation of Section~\ref{subsec:lift-exponent-primes}.
The depth and all other terms are therefore independent of the Frobenius and conjugation labels. One power applies at every maximal ideal in the cover. Outside the weight range use the sum of the full Weyl exponents, with exponent zero above height. The small and excluded primes likewise use the full Weyl sum. At a remaining prime the contribution of each larger-weight index is the minimum of its full exponent and its integral coarse correction. These comparisons fix the
tuple at \(w\) before taking the least exponent and keep the
original weight record for the prime sum. In compatible Smith bases the two lattices to be compared are
\[
 \mathcal C_w=\bigoplus_i O_we_i,\qquad \mathcal D_w=\bigoplus_i O_w\varpi_w^{d_i}e_i\subseteq\mathcal C_w.
\]

Then
\[
     J_{j,w}=\max_i d_i,\qquad
     p^{B_j(p)}\mathcal C_w\subseteq\mathcal D_w
             \ \Longrightarrow\ d_i\le e_wB_j(p)\quad\text{for all }i.
\]
The same comparison applies if the surviving eigensystem
requires an extension of the coefficient field. Over the extension the exponents and
the absolute ramification index both multiply by its relative
ramification index.
\label{cor:lift-exponent-models-3}
\end{proof}


\subsubsection{Summing the weight exponents}\label{subsec:lift-exponent-primes}
\hypertarget{source-lift-exponent-primes}{}

The integral inclusion gives local exponents indexed by highest weights. The sum over primes depends on which differences between shifted weights can be divisible by the prime. We keep the head weights and the consecutive block distinct in making this count.

\begin{established}[Heights, cutoff and exceptional interval]\label{hyp:weyl-prime-range}
Use the weight record for one lift. In this subsection \(r\) is
its coefficient rank, \(k_c=k-1\) is the head size, \(t\) is its
consecutive block length, and \(S,\eta\) bound the widths of
the head weights and the displaced weights, respectively.
The shifted weights in that
record are the highest weights plus the usual staircase. Their
height is bounded by \(h_0=\lfloor h\rfloor\), where
\(h=b+S+\eta+t+1\). Write \(X'\) for the auxiliary cutoff,
\(b'=xb/X'\), \(f=2n\), and \(w''=b/(nX')\).

Let \(N=k+m(t+1)\) be the final rank, and let \(\ell_*\)
be the auxiliary threshold
\[
 \ell_*=\frac{mb}{nt^2(N+t)z^2+kn^2(S+t+z)} .
\]
Put \(K_*=\lfloor 2b/\ell_*\rfloor\). Define the small-prime
threshold \(y\) and the exceptional interval \(\mathcal I\) by
\[
 y=\frac{mbz}{k(m+t)\log(3h)},\qquad
 \mathcal I=\left[b-\frac{t-1}{2}-K_*\,,\,
       b+S+\eta+1+\frac{t-1}{2}+K_*\right].
\]
Uniformly beyond a fixed threshold one has \(z\ge4\),
\(y>8\), \(y\le b\), \(r+1<p\) and \(S+k<p\) for \(p>y\),
\(h\le2b\), \(h+r\le3b<y^2\), and \(\mathcal I\subset(0,y^2)\).
Furthermore \(k(S+1)\le2kS\le2w''\) since \(S\ge1\). \label{cor:lift-exponent-primes-1}
\end{established}
\begin{proof}
Here \(b=\exp(x^\epsilon)\), \(z=\log b\). The ranks \(k,m,t,r\) are polynomially bounded in \(x\), and \(S,\eta\le x^{O(1)}\exp(x^{\epsilon/4})=b^{o(1)}\) by the weights of Lemma~\ref{lem:window-parameter-choice}. Hence \(h=(1+o(1))b\), \(K_*=b^{o(1)}\) by its displayed denominator, and \(\log(3h)\ge z\). As \(m\le k(m+t)\), one has \(b^{1-o(1)}\le y\le b\). In particular \(y^2>3b\), \(r+1,S+k<y\), and both endpoints of \(\mathcal I\) equal \(b+o(b)\). Finally \(kS=b^{o(1)}\) and \(w''=b^{1-o(1)}\), which proves the last comparison.
\end{proof}

For an integral highest weight \(\lambda\) in rank \(d\), let
\(e_p(\lambda)\) be the least exponent carrying the dual Weyl
lattice into the Weyl lattice over \(\Z_{(p)}\). Determinant
translation of both lattices leaves it fixed.

\begin{lemma}[Cutting a Weyl lattice at a consecutive block]
\label{lem:consecutive-weyl-cut}
The exponent satisfies the following estimate at a consecutive cut into ranks
\(a\) and \(d-a\). If the full shifted weights lie in an interval
of integer width \(W\), then
\begin{equation}\label{eq:lift-weyl-cut}
 e_p(\lambda)\le e_p(\lambda_{<a})+e_p(\lambda_{\ge a})
                   +a(d-a)\lfloor\log_p W\rfloor .
\end{equation}
Use zero for the logarithmic term at \(W=0\). Here the
restricted weights use the staircase in their own ranks.
Thus restriction changes the shift by a constant on each block
and preserves its width. \end{lemma}

\begin{proof}
We require a rational operator for the Levi blocks that realizes the two integral inclusions with scalars \(c_1,c_2\). In the adjoint action on endomorphisms,
apply the successive rank-one projections for the cross-block
pairs. Grid invariance gives an endomorphism of the full
representation. Its scalar on the highest weight vector is the
inverse of the product of the cross-block shifted differences.
Schur's lemma gives that same scalar on the module. In the
integral grid comparison the beta denominator for each pair
is bounded by the denominators of
\(i!j!/(i+j+1)!\) for \(i+j+1\le W\). Put \(m=p^a\).

At this power the factorial valuation contributes at most one, since
\[
 \left\lfloor\frac{i+j+1}{m}\right\rfloor-\left\lfloor\frac{i}{m}\right\rfloor-\left\lfloor\frac{j}{m}\right\rfloor=\left\lfloor\frac{(i\bmod m)+(j\bmod m)+1}{m}\right\rfloor\le1 .
\] Hence a factor
\(p^{\lfloor\log_p W\rfloor}\) per pair, times the two block
inclusion scalars, clears the full inclusion once the grid is
multiplied by its highest-weight scalar denominator. This gives the grid identities for the rank-one line, its scalar and the integral denominator. We use those grid identities
on the Weyl lattices with the stated width condition.
\label{cor:lift-exponent-primes-2}

\end{proof}

\begin{proposition}[Sum of least integral exponents on the selected record]
\label{prop:selected-weyl-prime-sum}
With the heights, cutoff and exceptional interval of
Lemma~\ref{hyp:weyl-prime-range}, for
each selected lift the original rational-prime contributions
\(B_j(p)\) in \eqref{eq:division-least-sum} have
\[
 \sum_p B_j(p)\log p\le30|T_2|mbz^2+2mb.
\]
The estimate also holds for a restricted set of rational primes
satisfying Lemma~\ref{hyp:weyl-prime-range}. The constants and threshold are
uniform over the selected lifts and coefficient fields.
\end{proposition}

\begin{proof}
\proofstep{The width estimate and the small primes.}
Rank zero and rank one have exponent zero. Iterating
\eqref{eq:lift-weyl-cut} with one rank-one block gives
\(e_p(\lambda)\le\binom d2\lfloor\log_p W\rfloor\), and hence
zero for \(p>W\). At an index with a smaller weight, the width satisfies \(W_\tau\le2b'\).
If \(W_\tau>1\), the bound
\(\pi(Y)<1.26Y/\log Y\) for \(Y>1\), used here with
\(Y=W_\tau\), follows from Rosser and Schoenfeld's explicit
estimate \cite{rosser-schoenfeld-explicit}. This gives
\[
 \sum_p e_p(\lambda_\tau)\log p
 \le \binom r2\log W_\tau\,\pi(W_\tau)
 \le 1.3r^2 b'.
\]

If \(W_\tau\le1\), this sum is zero. At an index with a larger weight the
crude pointwise bound is \(\binom r2\log h\) for the
exponent times \(\log p\). For \(p\ge2\eta+t\), cut off the
head of size \(k_c\). Its weights lie in an interval of width
at most \(S\). The
remaining weights have width strictly less than \(p\) under
the weight shape conditions. The lower block therefore has
exponent zero, and \eqref{eq:lift-weyl-cut} gives the bound
\[
 e_p(\lambda_\tau)\log p\le
 k(m+t)\log(3h)+k^2\log(S+k).
\]

In particular, the head by itself has zero exponent when
\(p>S+k\). At \(p\le y\), the displayed bound times \(p\)
is at most \(2mbz\) since \(k\le m+t\), \(S+k\le3h\),
and \(y\,k(m+t)\log(3h)=mbz\). In the other case
\(p<2\eta+t\), the size conditions \(t\le\eta\), \(r\le2t\)
and \(\eta\le b/(t^3 f k^3)\) give the same bound from
\(e_p(\lambda_\tau)\log p\le4t^2 z\) and
\(m f t k^3\ge6\). For an index with a larger weight the small-prime contribution is at most \(2mbz^2\), since
\[
 \sum_{p\le y}\frac1p\le\sum_{j=2}^{\lfloor y\rfloor}\frac1j\le \log y\le z .
\]

\proofstep{The exceptional primes above the cutoff.}
Above \(y\), an exceptional prime divides an integer of
\(\mathcal I\). Each such integer has at most one distinct
prime factor above \(y\), as all these integers are positive
and less than \(y^2\). Hence the number of excluded primes
above \(y\) is at most \(S+\eta+t+2K_*+1\), plus the
number of specified primes. The conductor of the auxiliary field factors that
are at most \(y\) are already charged in the small-prime range.
Indeed the additional small factors divide an integer at most
\(n^2\le kt\), and \(kt\le b/(4kt)\le y\) by the large
parameter conditions.

The four listed primes \(q,d,p_{\rm de},
p'_{\rm de}\) suffice in the specified set after this removal.
Allowing five gives an exceptional cost at most
\begin{equation}\label{eq:exponent-exceptional-prime-cost}
 (S+\eta+t+2K_*+6) f r^2 z,
\end{equation}
because each of the \(f\) label contributions satisfies
\[
 e_p(\lambda_\tau)\log p
 \le \binom{r}{2}\log h\le r^2 z
\]
whenever its width exceeds one. For a smaller weight use its
width, at most \(2b'\) and at most \(h\), in place of \(h\);
widths at most one have zero exponent.
The vanishing above height holds on the full record because the
smaller widths are at most \(h\) as well.
\label{cor:lift-exponent-primes-3}

\proofstep{The restricted exponent at the remaining primes.}
At primes \(y<p\le h\) outside the excluded interval and the
specified set, use the local weight-elimination formula with
the depth estimate of Lemma~\ref{lem:selected-pieri-depth}. Let \(j_\tau,q_\tau\ge0\), \(j_\tau+q_\tau\le m\),
be its two counts outside the head and consecutive block. At an index with a larger weight, the head interval \([x_-,x_+]\) and the top weight \(B_\tau\) of the consecutive block satisfy
\[
 0\le B_\tau<x_-\le x_+\le h,\qquad x_+-x_-\le S.
\]
Put \(d_\tau=x_--B_\tau\).

The head exponent in this formula is zero. The correction counts from the windows
around the consecutive blocks of radius \(\lfloor h/(p-1)\rfloor\)
are zero for the original and the Frobenius-labelled weights.
In the signed coordinates of the shape condition, center the
comparison at the midpoint of the consecutive block. Write \(c\) for this midpoint, \(z_0\) for a point of the
enlarged consecutive window, and \(H_p=\lfloor h/(p-1)\rfloor\).
The window and the signed displacement bounds give, for a displaced
weight \(u\),
\begin{equation}\label{eq:weyl-short-window-distance}
 |z_0-c|\le (t-1)/2+H_p,\qquad
 0<|u-z_0|\le\eta+t+H_p<p .
\end{equation}
The strict lower bound uses the lower displacement in the weight
record. Thus no displaced weight is congruent to \(z_0\) modulo
\(p\).

For the head weights there are two further tests. At an
index with a larger weight the signed difference, with the sign
prescribed by the shape, belongs to \(\mathcal I\).
At a top or an index with a smaller weight the absolute
difference lies in \((0,p)\). Summarizing the exclusions,
\[
 \begin{aligned}
 d\in\mathcal I\cap\mathbb Z&\quad\Longrightarrow\quad p\nmid d
       &&\text{by the choice of }p,\\
 0<|d|<p&\quad\Longrightarrow\quad p\nmid d
       &&\text{by size}.
 \end{aligned}
\]
Reflection preserves the absolute differences in \eqref{eq:weyl-short-window-distance} and the nondivisibility tests, so the reflected shape has the same vanishing correction counts.
Thus the remaining depth in the real bound for
this formula is \(q_\tau h\). Let \(I_p\) indicate whether
\(p\mid u-B_\tau\) for an integer \(u\in[x_-,x_+]\).

Write \(r_p\) for the least positive residue of \(d_\tau\) modulo
\(p\). The term \(G_p\) is \(\min(k,t,q_\tau h/r_p)\) if
\(r_p<d_\tau\), and zero otherwise. The larger-weight
part of the rational-prime summand, before multiplying by
\(\log p\), is bounded by
\[
 (q_\tau+2j_\tau)\left(\frac{h+r}{p-1}+6\right)
   +\frac{kS}{p}+k I_p+G_p
   +\frac{q_\tau h}{p}(1+\log(h/p)).
\]
Two elementary sums give
\[
 \sum_{y<p\le h} I_p\log p\le(S+1)\log h,\qquad
 \sum_{y<p\le h} G_p\log p\le q_\tau h(1+\log h)\log h.
\]

\proofstep{Summing the two residue terms.}
For the first, map each prime counted by the indicator to one
integer \(u\in[x_-,x_+]\) such that \(p\mid u-B_\tau\). Two distinct primes cannot give the same \(u\). Otherwise
\[
 y^2<pp'\le u-B_\tau\le h<y^2 .
\] The interval contains at most \(S+1\) such integers.
For the second, every prime with \(G_p\ne0\) divides the positive
integer \(d_\tau-r_p\le h\). Thus \(p\mapsto r_p\) is injective
by the same product inequality. Sum \(q_\tau h/r\) over
\(1\le r\le\lfloor h\rfloor\), and bound \(\log p\) by \(\log h\).
\label{cor:lift-exponent-primes-4}
\proofstep{Elementary weighted prime sums.}
The other sums use \(\sum_{p\le h+r}\log p\le(h+r)\log4\)
and \(\sum_{p\le h+r}(\log p)/p\le\log(h+r)+\log4\).
At any integer \(Y\ge2\), sum \(Y/p\le\lfloor Y/p\rfloor+1\) with weights \(\log p\). The contributions of the two sums give
\[
 Y\sum_{p\le Y}\frac{\log p}{p}\le \sum_{p\le Y}\lfloor Y/p\rfloor\log p+\sum_{p\le Y}\log p\le\log(Y!)+Y\log4 .
\]
The floor terms are bounded by \(\log(Y!)\) and the unit terms
by \(Y\log4\). Division by \(Y\) and taking floors in the
cutoff give the bound for real cutoffs as well.

For \(p>y>8\) one can replace \(1/(p-1)\) by \(9/(8p)\).
Here \(h+r\le3b\), \(\log h\le z+0.7\), and
\(\log(h+r)+\log4\le z+2.49\). The two terms with factor
\(k\) contribute at most \(2k(S+1)(z+2)\) since
\(S(z+2.49)+(S+1)(z+0.7)\le2(S+1)(z+2)\).
The remaining four terms together cost at most \(mb\) times
\[
 \tfrac{27}{4}(z+2.49)+50.04
   +2(z+1.7)(z+0.7)+2(z+1.7)(z+2.49)\le18z^2.
\]
For the last inequality the difference has value and derivative
positive at \(4\) and has positive quadratic coefficient.
These bounds apply to every subset of the indicated primes.

\proofstep{Absorption of the three tails.}
After summing the residue terms, the remaining tails involve different sources of divisibility. Their bounds must all fit the exponent allowance for this single lift.

Let \(T_2'\) be the number of indices with larger weights in the record.
Adding the three prime ranges gives the term
\(30T_2' mbz^2\), with room from \(2+18\le30\). The remaining
terms are
\begin{equation}\label{eq:exponent-prime-tails}
 2T_2' k(S+1)(z+2)+1.3 f r^2 b'
       +(S+\eta+t+2K_*+6) f r^2 z .
\end{equation}

To bound them let \(A\ge1\) be the polynomial size majorant
for the selected parameters. The active size rows give
\[
 \begin{gathered}
 f r^2\le4A^3,\quad kS\le w'',\quad S+\eta\le b/X',\quad
 K_*+t\le A^4x^2(z^2+S),\\
 4(z+2)\le X',\quad 5.2A^3x\le X',\quad
 40 A^7x^2 z\le X',\quad 40 A^7x^2 z^3\le b .
 \end{gathered}
\]
The bound for \(K_*+t\) follows by expanding \(2b/\ell_*\)
and using \(n,k,t\le A\), \(N\le Ax\) and \(m\ge2\).
The last four inequalities hold eventually for the common
polynomial majorant at \(X'=\exp(x^{\epsilon/3})\),
\(z=x^\epsilon\), \(b=\exp z\). The first term of
\eqref{eq:exponent-prime-tails} is at most
\(8(z+2)b/X'\le2b\), using \(T_2'\le f=2n\).
The second is at most \(5.2A^3xb/X'\le b\).

For the third use
\(S+\eta+t+2K_*+6\le5A^4x^2(z^2+b/X')\).
After multiplication by \(fr^2z\) the two terms are each
at most \(b/2\). Adding these three remainders and using \(m\ge2\) gives
\[
 2b+b+2(b/2)=4b\le 2mb .
\] This proves the rational-prime bound in
\eqref{eq:division-least-sum} with \(T_2'\le|T_2|\).
It is uniform over the chosen weights, lifts, coefficient fields
and subsets of rational primes in those size rows.
\label{cor:lift-exponent-primes-5}
\end{proof}


\paragraph{A common class denominator.}
For each lift \(j\) and tuple \(t\), the integer \(D_{G,j,t}\)
used above is the order of the row level subgroup times the
least common multiple of its class stabilizer orders. Take
\(D_*=\operatorname{lcm}_{j,t}D_{G,j,t}\).
The selected row bundle identifies the level as a subgroup of
the common finite group \(G\). Its stabilizer orders divide
the positive integer \(D_{\Gamma,u_0}\) of the distinguished
word \(u_0\). These divisibilities hold for every tuple in the
bundle, including all \(\mu(w)\). It follows that
\[
                   D_*\mid |G|D_{\Gamma,u_0}.
\]
This argument uses a single common multiple even when the
number of tuples grows. In the active fixed-scale comparison
the value-lattice multiplier is \(s=1\).

The row finite group has
\(\log|G|\le 2fr^2\log(2g_0)\).
Indeed each of its residue factors is bounded by the order of
the corresponding unitary or general linear matrix group.
In the unitary case its order is
\[
 q^{r^2}\prod_{i=1}^r(1-(-q)^{-i})\le (2q)^{r^2}.
\]

The local degrees over \(g_0\) and the selected level size give
the stated group bound. The size conditions give \(2r^2\le b\)
and \(\log(2g_0)\le2z^2\), so \(\log|G|\le2\mathcal E\).
The stabilizer comparison has the bound
\[
 \log D_{\Gamma,u}\le C_\Gamma fN_m\log(fN_m)
\]
for every
word \(u\), where \(N_m\) is the final unitary rank. Here the
stabilizers act faithfully on rational spaces of dimension
at most \(2fN_m\) by restriction of scalars. Minkowski's
finite-subgroup bound gives a common multiple independent
of the lattice.

For rational dimension \(n\), its exponent
at a prime \(p\) is
\(\sum_{a\ge0}\lfloor n/(p^a(p-1))\rfloor\).
The logarithm of that common multiple is \(O(n\log(2n))\):
only \(p\le n+1\) occur, and summing the exponents is
bounded by \(n\sum_{p\le n+1} p\log(p)/(p-1)^2\).
Chebyshev's bound and partial summation give
\[
 \sum_{p\le u}\frac{\log p}{p}=\frac{1}{u}\sum_{p\le u}\log p+\int_2^u\frac{\sum_{p\le t}\log p}{t^2}\,dt=O(\log(2u))\quad(u=n+1).
\]
\cite{serre-finite-orders}
The conductor and rank size condition gives
\(1\le fN_m\le3\sqrt b\). In particular
\(\log(fN_m)\le3z\) and
\(fN_m\log(fN_m)\le9bz\le9\mathcal E\).
Thus on the selected row bundle
\begin{equation}\label{eq:division-class-bound}
                \log D_*\le(2+9C_\Gamma)\mathcal E .
\end{equation}
The coefficient field and the number of words, lifts and
tuples cause no additional multiplier in this bound.
\label{cor:division-length-9}
